\documentclass{article}

\usepackage{amsmath, amsthm, amssymb, amsfonts}
\usepackage{thmtools}
\usepackage{graphicx}
\usepackage{setspace}
\usepackage{geometry}
\usepackage{float}
\usepackage{hyperref}
\usepackage[utf8]{inputenc}
\usepackage[english]{babel}
\usepackage{framed}
\usepackage[dvipsnames]{xcolor}
\usepackage{tcolorbox}

\colorlet{LightGray}{White!90!Periwinkle}
\colorlet{LightOrange}{Orange!15}
\colorlet{LightGreen}{Green!15}

\newcommand{\HRule}[1]{\rule{\linewidth}{#1}}

\declaretheoremstyle[name=Theorem,]{thmsty}
\declaretheorem[style=thmsty,numberwithin=section]{theorem}
\tcolorboxenvironment{theorem}{colback=LightGray}

\declaretheoremstyle[name=Proposition,]{prosty}
\declaretheorem[style=prosty,numberlike=theorem]{proposition}
\tcolorboxenvironment{proposition}{colback=LightOrange}

\declaretheoremstyle[name=Principle,]{prcpsty}
\declaretheorem[style=prcpsty,numberlike=theorem]{principle}
\tcolorboxenvironment{principle}{colback=LightGreen}

\setstretch{1.2}
\geometry{
    textheight=9in,
    textwidth=5.5in,
    top=1in,
    headheight=12pt,
    headsep=25pt,
    footskip=30pt
}

% ------------------------------------------------------------------------------

\begin{document}

% ------------------------------------------------------------------------------
% Cover Page and ToC
% ------------------------------------------------------------------------------

\title{ \normalsize \textsc{}
		\\ [2.0cm]
		\HRule{1.5pt} \\
		\LARGE \textbf{\uppercase{ECE 6553 - Optimal Controls}
		\HRule{2.0pt} \\ [0.6cm] \LARGE{Spring 2025} \vspace*{10\baselineskip}}
		}
\date{}
\author{\textbf{Authors} \\ 
		Jackson Isenberg \\
		Erik Verriest}

\maketitle
\newpage

\tableofcontents
\newpage

% ------------------------------------------------------------------------------

\section{Reachability of Linear Systems}

\subsection{Reachability of an LTI System}

Consider the Linear Time-Invariant (LTI) finite system
\begin{align*}
    \dot{x} &= Ax + bu \\
    y &= cx
\end{align*}
with $A \in \mathbb{R}^{n \times n}$, $b \in \mathbb{R}^{n \times 1}$, $c \in \mathbb{R}^{1 \times n}$.
We define our system to be $\Sigma = (A, b, c)$. The reachability matrix $\mathcal{R}$ is
\[\mathcal{R}(A, b) = \left[\begin{matrix}
    b & Ab & \cdots & A^{n-1}b
\end{matrix}\right]\]
We determine reachability for the system by looking to see if $\det\mathcal{R}(A,b) \ne 0$.
In other words, if the reachability matrix is invertible (non-singular), then we have reachability.
Reachability means that given some initial state, the control system can drive it to
any other state within its state space.

\noindent
Say we are given $u(t)$ and $y(t)$ for $t > 0$. Can we find $x(0)$? The observability matrix
$\mathcal{O}$ is
\[\mathcal{O}(A, c) = \left[\begin{matrix}
    c \\
    cA \\
    \vdots \\
    cA^{n-1}
\end{matrix}\right]\]
The system is observable if $\det\mathcal{O}(A,c) \ne 0$, which means we would be able to determine
$x(0)$. observability means that the system's internal state can be fully understood
with only the external outputs.

\noindent
However, there is an easier way to determine reachability without computing the reachability matrix.
The Popov-Belerith-Hautus (PBH) test says that for the augmented matrix
\[\left[\begin{matrix}
    sI - A & \mid & b
\end{matrix}\right] \in \mathbb{R}^{n \times (n+1)}\]
that we can determine reachability if
\[\text{rank}\left[\begin{matrix}
    sI - A & \mid & b
\end{matrix}\right] = n, \forall s\]

\subsection{Fundamental Realization Theorem}

We can realize a transfer function for an LTI system $\Sigma = (A,b,c)$
if $(A,b)$ is reachable and $(A,c)$ is observable. This leads to the realization of
\[H(s) = \frac{b(s)}{a(s)} = \frac{c \text{Adi}(sI - A)b}{\det(sI-A)}\]
The adioint (or adiugate) matrix for a matrix $A$ is defined as the matrix that satisfies
\[A^{-1} = \frac{1}{\det(A)}\text{Adi}(A)\]
and is also the transpose of the cofactor matrix of $A$. The cofactor matrix $C$ of $A$ is defined
to have the following elements:
\[C_{ii} = (-1)^{i+i}\det(M_{ii})\]
where $M_{ii}$ is the minor matrix created by removing the $i$-th row and $i$-th column of $A$.
So then the adioint matrix would be
\[\text{Adi}(A)_{ii} = C_{ii} = (-1)^{i+i}\det(M_{ii})\]

\noindent
The \textit{minimal} realization is when $a(s)$ and $b(s)$ are coprime or relatively prime.

\subsection{Linear Time-Variant Systems}

We have looked at time-invariant systems, but what if the behavior of the system
also depends on time? Consider the scalar system
\begin{align*}
    \dot{x}(t) &= a(t)x(t), \quad x(0) = x_0
\end{align*}
We see that
\[\frac{\dot{x}(t)}{x(t)} = \frac{d}{dt}\ln x(t) = a(t)\]
\[\ln x(t) - \ln x(0) = \int_{0}^{t}a(\tau)d\tau\]
\[\ln x(t) = \ln x(0) + \int_{0}^{t}a(\tau)d\tau\]
\[x(t) = x(0)\exp\int_{0}^{t}a(\tau)d\tau\]
However, this is only true for the scalar case. For the vector case, the solution
is more complex:
\[x(t) = x_0 + x_0\int_{0}^{t}A(\tau)d\tau + x_0\int_{0}^{t}A(\tau)\int_{0}^{\tau}A(\tau_1)d\tau_1d\tau
+ \cdots \]
Then, for a non-zero initial value, we end up with:
\[x(t) = x(\tau) + x(\tau)\int_{\tau}^{t}A(\theta)d\theta + \cdots\]
\[= \phi(t, \tau)x(\tau)\]
where $\phi(t, \tau)$ is our transition matrix from time $\tau$ to time $t$.
Now, consider the system
\begin{align*}
    \dot{x}(t) &= A(t)x(t) + b(t)u(t)
\end{align*}
The solution is then
\[x(t) = \phi(t,\tau)x(\tau) + \int_{\tau}^{t}\phi(t,\theta)b(\theta)u(\theta)d\theta\]

\subsection{Instantaneous Reachability}

Consider the system
\[\dot{x}(t) = A(t)x(t) + b(t)\delta(t)\]
which has impulsive input. We know that $\delta(t) = H'(t)$, so
\begin{align*}
    \dot{x}(t) &= A(t)x(t) + b(t)\frac{d}{dt}H(t) \\
    &= A(t)x(t) + \frac{d}{dt}[b(t)H(t)] - H(t)\frac{d}{dt}b(t) \quad \text{(by product rule)}
\end{align*}
\begin{align*}
    \frac{d}{dt}[x(t)-b(t)H(t)] &= A(t)x(t) - H(t)\frac{d}{dt}b(t) \\
    &= A(t)[x(t) - b(t)H(t)] + H(t)[A(t)b(t)-\dot{b}(t)]
\end{align*}
This is of the form
\[\frac{d}{dt}z(t) = A(t)z(t) + \text{discontinuous function}\]
where $z$ is continuous.

\noindent
From further derivation, we can see that for the system
\[\dot{x}(t) = A(t)x(t) + b(t)\delta^{(k)}(t)\]
that
\[x(0+) = \left([A(t)-D]^k b(t)\right)_{t = 0} + x(0-)\]
where the operator $(\cdot)_{t=0}$ means evaluated at $t = 0$ and $D$ is the derivative
operator such that $Df = \dot{f}$. Consider now the input
\[u(t) = \sum_{i=0}^{N-1}\alpha_i \delta^{(i)},\]
we get that
\[x(0+) = x(0-) + \left[b, (A-D)b, (A-D)^2 b, \ldots\right]\left[\begin{matrix}
    \alpha_0 \\
    \alpha_1 \\
    \alpha_2 \\
    \vdots
\end{matrix}\right]\]
So, the infinite reachability matrix is
\[\mathcal{R}(A,b) = \left[b, (A-D)b, (A-D)^2 b, \ldots\right]\]
Note that the entries are functions dependent on time $t$, so the time-variant
reachability also depends on $t$.

\subsection{Finite Time Reachability}

For $t \rightarrow T$,
\[x_f = x(T) = \phi(T, \tau)x(\tau) + \int_{\tau}^{T}\phi(T,\theta)b(\theta)u(\theta)d\theta\]
Then, let $u(\theta)$ be of the form
\[u(\theta) = b(\theta)^\top \phi^\top(T,\theta)q\]
where $q$ is some unknown constant vector. Our solution becomes
\[x(T) = \phi(T, \tau)x(\tau) +
q\int_{\tau}^{T}\phi(T,\theta)b(\theta)b(\theta)^\top\phi^\top(T,\theta)d\theta\]
We call this the reachability Gramian, defined as
\[\mathcal{R}_{(t_0, t_f)}(A,b) = \int_{t_0}^{t_f}\phi(t_f,t)b(t)b^\top(t)\phi^\top(t_f,t)dt\]
The optimal input $u^*(t)$ is the input that will drive our state from $(x_0,t_0) \rightarrow (x_f,t_f)$
where $t_f > t_0$. We define it as such:
\[u^*(t) = b^\top(t)\phi^\top(t_f,t)\left[\mathcal{R}_{(t_0, t_f)}(A,b)\right]^{-1}
(x_f - \phi(t_f,t_0)x_0)\]
A unique property of this optimal input is that it is the input with minimal energy:
\[u^*(t) = \min_{u}\mathcal{E}_u\]
\[\mathcal{E}_u = \frac{1}{2}\int_{t_0}^{t_f}u^2(t)dt\]

\section{Linearization of a Nonlinear System}

Consider the system
\begin{align*}
    \dot{x} &= f(x, u) \\
    y &= h(x, u)
\end{align*}
\[x \in \mathbb{R}^{n}, f = \left[\begin{matrix}
    f_1 \\
    \vdots \\
    f_n
\end{matrix}\right]\]
We say that $x_e$ is an equilibrium point if $f(x_e, 0) = 0$. Stability of nonlinear
systems is defined by the stability of the equilibrium points, not of the system itself. But how
do we find if an equilibrium point is stable or not (i.e., look at the behavior of $\dot{x}$
around $x_e$)?

\noindent
We perturb the equilibrium point by a small amount $\tilde{x}$: $x = x_e + \tilde{x}$. Then,
\[\dot{x} = \dot{x}_e + \dot{\tilde{x}} = f(x_e + \tilde{x}, 0)\]
Since $\dot{x}_e$ is a constant, the derivative is 0. So,
\[\dot{\tilde{x}} = f(x_e + \tilde{x}, 0) = f(x_e, 0) +
\left.\frac{\partial f}{\partial x}\right|_{x_e, 0}\tilde{x} + \text{higher order terms}\]
where
\[\frac{\partial f}{\partial x} = \left[\begin{matrix}
    \frac{\partial f_1}{\partial x_1} & \cdots & \frac{\partial f_1}{\partial x_n} \\[2mm]
    \vdots & \ddots & \vdots \\[2mm]
    \frac{\partial f_n}{\partial x_1} & \cdots & \frac{\partial f_n}{\partial x_n}
\end{matrix}\right]\]
is the Jacobian. Since $\tilde{x}$ is small, we can rewrite our system as
\[\dot{\xi} = \left.\frac{\partial f}{\partial x}\right|_{(x_e, 0)}\xi\]
with $\xi \approx \tilde{x}$. We see that this is a linear time-invariant system. This is an example
of linearizing a nonlinear system.

\noindent
Let us look at the behavior of the system around the initial values $(x_0, u_0)$.
We perturb $x_0$ and $u_0$:
\begin{align*}
    x &= x_0 + \tilde{x} \\
    u &= u_0 + \tilde{u}
\end{align*}
So,
\[\dot{x} = \dot{x}_0 + \dot{tilde{x}} = f(x_0 + \tilde{x}, u_0 + \tilde{u})\]
\[\dot{\tilde{x}} = f(x_0, u_0) + \left.\frac{\partial f}{\partial x}\right|_{(x_0, u_0)}\tilde{x}
+ \left.\frac{\partial f}{\partial u}\right|_{(x_0, u_0)}\tilde{u} + \text{h.o.t.}\]
Then our linearized system with $\xi \approx \tilde{x}$ is
\[\dot{\xi} = \left.\frac{\partial f}{\partial x}\right|_{(x_0(t), u_0(t))}\xi
+ \left.\frac{\partial f}{\partial u}\right|_{(x_0(t), u_0(t))}\tilde{u}\]
which is a linear time-variant system.

\subsection{Stability of the Equilibrium}

Consider the equilibrium point $(x_e, u_e)$. For input $u(t) = u_e$, we say that
the equilibrium point is stable if the solution $x(t)$ converges to $x_e$. Let the Jacobian
$\frac{\partial f}{\partial x}$ evaluated at $(x_e, u_e)$ be the matrix $A$. The linearized system
is $\dot{\xi} = A\xi$.

\noindent
If $\text{Re}\lambda(A) < 0$, then $(x_e, u_e)$ is locally asymptotically stable.

\noindent
If there is an eigenvalue of $A$ with a positive real part, then the equilibrium is locally unstable.

\noindent
If an eigenvalue of $A$ is purely imaginary, then no conclusion can be drawn.

\subsubsection{Hartman-Grossman Theorem}

The equilibrium $(x_e, u_e)$ for $\dot{x} = f(x, u)$ is said to be \textit{hyperbolic}
if $\text{Re}\lambda(A) \ne 0$ where $A = \frac{\partial f}{\partial x}$ evaluated at $(x_e, u_e)$.

\begin{theorem}
    If $(x_e, u_e)$ is a hyperbolic equilibrium, then the global stability of the linearized model
    implies the local stability of the equilibrium for the nonlinear system.
\end{theorem}

\section{Parameter Optimization}

Given \[L(x) \in \mathbb{R}, \quad x \in \mathbb{R}^n\] we want to find
the optimal parameter $x^*$ such that $L(x^*)$ is an extremum. The things we have to investigate
to do this are:
\begin{enumerate}
    \item Existence: ``Does an extremum even exist?''
    \item Uniqueness: Guarantee that different methods converge to the same solution
    \item Sufficient conditions: If they are satisfied, we can say we have an extremum
    \item Necessary conditions: Must be satisfied but do not guarantee an extremum
    \item Computation methods
\end{enumerate}

\subsection{Computational Methods}

\subsubsection{Geometric Methods}

\subsubsection{Algebraic Methods}

The algebraic-geometric mean inequality says that for $x_i > 0$ for $i = 1,\ldots,n$,
the arithmetic mean (AM) is
\[\text{AM}\{x_i\} = \frac{1}{n}\sum_{i=1}^{n}x_i\]
and the geometric mean (GM) is
\[\text{GM}\{x_i\} = \left(\prod_{i=1}^{n}x_i\right)^\frac{1}{n}\]
which follow the following inequality:
\[\text{AM}\{x_i\} \ge \text{GM}\{x_i\}\]
The quadratic mean (QM) is
\[\text{QM}\{x_i\} = \sqrt{\frac{1}{n}\sum_{i=1}^{n}x_i^2}\]
and then $\text{AM}\{x_i\} \le \text{QM}\{x_i\}$.
Finally, the harmonic mean (HM) is
\[\text{HM}\{x_i\} = \left[\frac{1}{n}\sum_{i=1}^{n}\frac{1}{x_i}\right]^{-1}\]
It holds that
\[0 < \text{HM} \le \text{GM} \le \text{AM} \le \text{QM}\]
\textbf{Matrix Inequalities:}

\noindent
Let $A \in \mathbb{R}^{n\times n}$ be a matrix. Then let
\[A_s = \frac{A+A^\top}{2}\]
\[A_A = \frac{A-A^\top}{2}\]
then, $A_s^\top = A_s$ and $A_A^\top = -A_A$.

\noindent
Now we look at the quadratic form. Let $P$ be a symmetric matrix.
\[x^\top Px, \quad x \in \mathbb{R}^{n}\]
Note the properties:
\[x^\top Ax = x^\top A_s x\]
\[x^\top A_A x = 0\]
Consider:
\[\min_{\|x\|=1}x^\top Px =\text{ ?}\]
Recall that
\[P = V\Lambda V^{-1} = V\Lambda V^\top\]
Then,
\[x^\top Px = x^\top V\Lambda V^\top x\]
Let $y = V^\top x$:
\[x^\top V\Lambda V^\top x = y^\top \Lambda y = \sum_{i=1}^{n}y_i^2\lambda_i\]
Note that
\[\|y\|^2 = y^\top y = x^\top V V^\top x = x^\top x = \|x\|^2\]
so length does not change. Therefore if $\|x\| = 1$, then $\|y\| = 1$.
So,
\[\min_{\|x\|=1}x^\top Px = \min_{\|y\|=1}y^\top \Lambda y =\text{ minimum eigenvalue of }P\]
\begin{theorem}
    Courant-Fisher Theorem\\[2mm]
    Consider that $x \perp \text{span}\{w_1,\ldots,w_{n-k}\}$, then
    \[\max_{\|x\|=1}x^\top Px = \lambda_k\]
\end{theorem}

\subsubsection{Analytic Methods}

\begin{theorem}
    Extreme Value Theorem\\[2mm]
    If $f$ is continuous with mapping from a compact (closed and bounded) set to $\mathbb{R}$, then
    $f$ will always have a maximum and a minimum.
\end{theorem}
\noindent
A counterexample is seen by considering $f(t) = e^t$ in the domain $(a, b)$
where $\infty < a < b < \infty$. Then $e^a < e^t, \forall t \in (a, b)$.
We say that $e^a$ is the infimum of $\{e^t\}$ in $(a,b)$.

\noindent
\textbf{Weistrass's Theorem:}

\noindent
For the 1-D case:
\begin{enumerate}
    \item Consider the necessary conditions for extremum in $[a, b]$
    \item The sufficient conditions for a local minimum at $x^* \in [a, b]$ are:
    \[f'(x^*) = 0\]
    \[f''(x^*) > 0\text{ if }x^*\text{ is in the interior }(a, b)\]
    \[\text{if }x^* = a \rightarrow f'(x^*) \ge 0, f''(x^*) > 0\]
    \item Existence: we see that if $f$ is continuous in $[a,b]$, then there must be a minimum
    \item Unique if $f$ is strictly convex, it as a minimum at a unique point in $[a,b]$
\end{enumerate}
\textbf{Derivatives with respect to Matrices:}

\noindent
Consider
\[L : \mathbb{R}^{n\times m} \rightarrow \mathbb{R}\]
\[L(X) = \text{Tr}(AX)\]
where $\text{Tr}$ is the trace function:
\[\text{Tr}(M) = \sum_{i=1}^{n}M_{ii}\]
with properties
\[\text{Tr}(M^\top) = \text{Tr}(M)\]
\[\text{Tr}(ABC) = \text{Tr}(BCA) = \text{Tr}(CAB)\]
Then,
\[\text{Tr}(AX) = \sum_{i=1}^{n}(AX)_{ii} = \sum_{i=1}^{m}\sum_{i=1}^{n}A_{ii}X_{ii}\]
The derivative with respect to $X$ is a matrix with elements
\[\left(\frac{\partial L}{\partial X}\right)_{ii} = \frac{\partial L}{\partial X_{ii}}\]
So,
\begin{align*}
    \left(\frac{\partial}{\partial X}\text{Tr}(AX)\right)_{ii} &=
    \frac{\partial}{\partial X_{ii}}\text{Tr}(AX) \\
    &= \frac{\partial}{\partial X_{ii}}\sum_{\alpha=1}^{m}
    \sum_{\beta=1}^{n}A_{\alpha\beta}X_{\beta\alpha} \\
    &= \sum_{\alpha=1}^{m}\sum_{\beta=1}^{n}A_{\alpha\beta}
    \frac{\partial X_{\beta\alpha}}{\partial X_{ii}} \\
\end{align*}
We then introduce the Kronecker Delta:
\[\delta_{\alpha\beta} = \begin{cases}
    1, & \alpha = \beta \\
    0, & \text{otherwise}
\end{cases}\]
Then,
\begin{align*}
    \left(\frac{\partial}{\partial X}\text{Tr}(AX)\right)_{ii} &=
    \sum_{\alpha=1}^{m}\sum_{\beta=1}^{n}A_{\alpha\beta}\delta_{\beta i}\delta_{\alpha i} \\
    &= \sum_{\beta=1}^{n}A_{i\beta}\delta_{b\beta i} \\
    &= A_{ii}
\end{align*}

\subsubsection{Numerical Methods}

\subsection{Parameter Optimization with Equality Constraints}

We define a problem with equality constraints to be one where
we want to find an extremum for $L(x)$ with constraint $f(x) = 0$. For example,
say there is a constraint that $x_1 = x_2$. Then, $f(x) = x_1 - x_2 = 0$ is our constraint.

\noindent
Let $L : \mathbb{R}^n \rightarrow \mathbb{R}$ be the function to be minimized.
We want to find $x^*$ such that $L(x^*)$ is minimal and $f(x^*) = 0$ with
$f : \mathbb{R}^n \rightarrow \mathbb{R}^m$ and $m < n$. We define $x$ to be
\[x = \left[\begin{matrix}
    x_1 \\ \vdots \\ x_n
\end{matrix}\right] = \left[\begin{matrix}
    y_1 \\ \vdots \\ y_m \\ u_1 \\ \vdots \\ u_{n-m}
\end{matrix}\right]\]
Then,
\[dL(x) = \partial L(y,u) = \frac{\partial L}{\partial y}dy + \frac{\partial L}{\partial u}du\]
\[= \left(\frac{\partial L}{\partial u}-\frac{\partial L}{\partial y}
\left(\frac{\partial f}{\partial y}\right)^{-1}\frac{\partial f}{\partial u}\right)du\]
We want this to be 0. We define
\[\lambda^\top = -\frac{\partial L}{\partial y}\left(\frac{\partial f}{\partial y}\right)^{-1}\]
These are the Lagrange multipliers. Our optimization goal becomes
\[\frac{\partial L}{\partial u}+\lambda^\top\frac{\partial f}{\partial u} = 0\]
We define the Hamiltonian $H$ to be
\[H = L + \lambda^T f\]
so that the optimization target (sufficient conditions) is
\[\frac{\partial H}{\partial u} = 0, \frac{\partial H}{\partial \lambda} = 0\]
With the necessary conditions
\[\frac{\partial H}{\partial x} = 0 \rightarrow \frac{\partial H}{\partial y} = 0,
\frac{\partial H}{\partial u} = 0\]
Example Problem:

\noindent
Let $L(x) = x_1^2 - x_2$ with the constraint $x_1 = x_2 + 1$.

\noindent
Then, $f(x) = x_1 - x_2 - 1$. So, our Hamiltonian is
\[H = L + \lambda^\top f = x_1^2 - x_2 + \lambda(x_1-x_2-1)\]
The necessary conditions are
\[\frac{\partial H}{\partial x_1} = 2x_1 + \lambda = 0 \rightarrow x_1 = -\frac{\lambda}{2}\]
\[\frac{\partial H}{\partial x_2} = -1 - \lambda = 0 \rightarrow \lambda = -1\]
So, $x_1 = \frac{1}{2}$. Then, we look at our constraint:
\[\frac{\partial H}{\partial \lambda} = x_1 - x_2 - 1 = 0 \rightarrow x_2 = x_1 - 1\]
So, $x_2 = -\frac{1}{2}$. Then, our optimal solution is
\[x^* = \left(\frac{1}{2},-\frac{1}{2}\right)\]

\subsection{Parameter Optimization with Inequality Constraints}

Instead of $f(x) = 0$, our problem becomes finding an extremum $x^*$ of $L(x)$
such that $f(x^*) \le 0$. Our Hamiltonian is
\[H = L + \mu^\top f\]
where $\mu \ge 0$. For this, we have to note the idea of inactive and active
constraints. A constraint is inactive if the extremum of the function already
satisfies $f(x) \le 0$. A constraint is active if the extremum lies in $f(x) > 0$
so we have to find the next smallest/largest value that does satisfy $f(x) \le 0$.
When solving for a parameter optimization problem like this, we look at each possible combination
of inactive and active constraints. For example, if there are two constraints,
we look at the cases where (a) no constraints are active, (b) the first constraint is active,
(c) the second constraint is active, (d) both constraints are active.

\noindent
Example Problem:

\noindent
Let $L(x) = x_1^2 - x_2$ with the constraints $x_2 \le x_1$ and $x_1 \ge 0$.

\noindent
We will define our constraint function $f$ to be
\[f(x) = \left(\begin{matrix}
    x_2 - x_1 \\
    -x_1
\end{matrix}\right) \le 0\]. So, our Hamiltonian is
\[H = L + \mu^\top f = x_1^2 - x_2 + \mu_1(x_2-x_1) + \mu_2(-x_1)
= x_1^2 - x_2 + \mu_1(x_2-x_1) - \mu_2x_1\]
The necessary conditions are
\[\frac{\partial H}{\partial x_1} = 2x_1 - \mu_1 - \mu_2 = 0 \rightarrow x_1 = \frac{\mu_1+\mu_2}{2}\]
\[\frac{\partial H}{\partial x_2} = -1 + \mu_1 = 0 \rightarrow \mu_1 = 1\]
We see that
\[x_1 = \frac{1+\mu_2}{2}\]
We then look at our constraints.
\begin{enumerate}
    \item Assume no constraints are active ($\mu_1 = \mu_2 = 0$)

    \noindent
    We know this cannot be the case because $\mu_1 = 1$ must be held. So, the first constraint
    needs to be considered.
    \item Assume the first constraint is active ($\mu_2 = 0$)

    \noindent
    Then, $x_1 = \frac{1}{2}$. We then look at our condition:
    \[\frac{\partial H}{\partial \mu_1} = x_1 - x_2 = 0\]
    So, $x_1 = x_2$ and therefore our optimal solution is
    \[x^* = \left(\frac{1}{2},\frac{1}{2}\right)\]
    which we see has $L(x_1,x_2) = \frac{1}{4} - \frac{1}{2} = -\frac{1}{4}$.
    \item Assume the second constraint is active ($\mu_1 = 0$)

    \noindent
    We know this cannot be the case because $\mu_1 = 1$ must be held. So, the first constraint
    needs to be considered.
    \item Assume both constraints are active.
    
    \noindent
    Since both constraints are active, then
    \[\frac{\partial H}{\partial \mu_1} = x_1 - x_2 = 0 \rightarrow x_1 = x_2\]
    \[\frac{\partial H}{\partial \mu_2} = -x_1 = 0 \rightarrow x_1 = 0\]
    So, $x_1 = x_2 = 0$ is the optimal solution here, with $L(0,0) = 0$.
\end{enumerate}
Knowing our two optimal solutions for the different cases, we see that $(1/2, 1/2)$
produces a more minimal point than $(0,0)$.

\section{Introduction to Optimal Control}

\subsection{Discrete Optimal Control}

Consider the discrete system
\[x_{k+1} = f(x_k, u_k), \quad x_0\text{ given }\]
with $k = 1,\ldots,N$. The optimization problem then is to find
\[\min_{\{u_k\}} \phi(x_N) + \sum_{k=0}^{N-1}L(x_k,u_k) = J(u_0,\ldots,u_N)\]
The direct solution is
\[x_k = X_k(u_0, u_1, \ldots, u_{k-1})\]

\subsubsection{One-Stage Optimal Control}

\[x_1 = f(x_0,u_0),\quad x_0\text{ given}\]
\[\min_{u_0}\left[\phi(x_1) + L(x_0,u_0)\right] = \min_{u_0,x_1}J(u_0,x_1)\]
Re-organize the Lagrange multiplier:
\[\bar{J} = J + \lambda_1^\top(f(x_0,u_0)-x_1)\]
which works because $x_1 = f(x_0,u_0)$, so this should evaluate to $\bar{J} = J$ optimally.
\[\bar{J} = \phi(x_1) + L(x_0,u_0) + \lambda_1^\top f(x_0,u_0) - \lambda_1^\top x_1\]
We notice that there is a Hamiltonian there:
\[H(x_0,u_0,\lambda_1) = L(x_0,u_0) + \lambda_1^\top f(x_0,u_0)\]
We then assume that $u_0^*$ is the optimum which produces the corresponding state update $x_1^*$.
Perturbing the system, we end up with
\[dJ = \left.\frac{\partial H}{\partial u}\right|_{x_1^*,u_0^*}du_0\]
Since the necessary condition for optimality is that the perturbation $dJ$ is zero, we
see then that it must hold that
\[\frac{\partial H}{\partial u} = 0\]

\subsubsection{Multi-Stage Optimal Control}

Consider the system
\[x_{k+1} = f_k(x_k, u_k), \quad x_0\text{ given }\]
with $k = 1,\ldots,N$. The performance index ($J$) is
\[J = \phi(x_N) + \sum_{k=0}^{N-1}L_k(x_k,u_k)\]
Following the same steps as with the one-stage optimal control, we end up with the perturbation
\[dJ = \sum_{k=0}^{N-1}\frac{\partial H_k}{\partial u_k}du_k\]
which should be 0. Which means that
\[\frac{\partial H_k}{\partial u_k} = 0\]
We also have the Lagrange multipliers
\[\lambda_k = \left(\frac{\partial H_k}{\partial x_k}\right)^\top,\quad
\lambda_N = \left(\frac{\partial\phi}{\partial x}\right)^\top\]
Note that we define the Hamiltonian to be
\[H_k = L_k + \lambda_{k+1}^\top f_k\]
Our system then has the costate equation of
\[\lambda_k^\top = \frac{\partial L_k}{\partial x_k} + \lambda_{k+1}^\top\frac{\partial f}{\partial x_k}\]
This means our state and costate equations are coupled in such a way that we would call this a
Two Point Boundary Value Problem. These are not relatively easy to solve.

\subsection{Calculus of Variations}

We cannot use this concept for systems in continuous time, so we introduce the concept of
calculus of variations. Let there be a performance index
\[J = \int_{x_0}^{x_1}F(x,y,y')dx\]
Why have the derivative $y'$ in the function? Without it, finding the optimum becomes a
much simpler problem. However, if we add constraints, then the optimal solution falls apart.
We need the derivative of $y$ so that we can guarantee that the boundary conditions are satisfied.

\noindent
We introduce an arbitrary function $\eta(x)$ which is continuous and differentiable in $[x_0, x_1]$.
Suppose then that $y^*(x)$ is the optimal solution. We perturb as such:
\[y(x) = y^*(x) + \epsilon\eta(x)\]
where $\epsilon$ is positive but goes to 0 (arbitrarily small).
Then the derivative is
\[y'(x) = y^{*'}(x) + \epsilon\eta'(x)\]
Our performance index (depending on $\epsilon$) is
\[J(\epsilon) = \int_{x_0}^{x_1}F(x,y^*(x)+\epsilon\eta(x),y^*(x')+\epsilon\eta'(x))\]
Now, let $\eta(x_0) = \eta(x_1) = 0$. Doing a Taylor expansion, we end up with the equation
\[\frac{\partial F}{\partial y} - \frac{d}{dx}\frac{\partial F}{\partial y'} = 0\]
which is the condition for optimality for $y^*(x)$. This is called the Euler-Lagrange equation.
\begin{theorem}
    Fundamental Lemma of Calculus of Variations\\[2mm]
    If $\varphi(t)$ is continuous in $[t_0,t_1]$ and if for all $\eta(t) \in C^1$ such that
    $\eta(t_0) = \eta(t_1) = 0$
    \[\int_{t_0}^{t_1}\eta(t)\varphi(t)dt = 0,\]
    then it must hold that
    \[\varphi(t) \equiv 0.\] 
\end{theorem}

\subsection{Optimal Controls with No Constraints}

Let our control problem be
\[\arg\min_{u(\cdot)}\int_{0}^{T}L(x(t),u(t),t)dt + \phi(x(T))\]
such that
\[\dot{x}(t) = f(x(t),u(t),t)\]
with $T < \infty$ being fixed. We can rewrite $L$ to be
\[L(x,u,t) = L(x,u,t) + \lambda^\top(t)(f(x,u,t)-\dot{x}) = H(x,u,t,\lambda)-\lambda\dot{x}\]
Then our performance index is
\[J = \phi(x(T)) + \int_{0}^{T}[H(x,u,t,\lambda)-\lambda\dot{x}]dt\]
which then can be rewritten as
\[J = \int_{0}^{t}[H + \dot{\lambda}^\top x]dt + \phi(x(T)) - \lambda^\top(T)x(T) + \lambda^\top(0)x(0)\]
Through this we end up with
\[\dot{\lambda} = -\left(\frac{\partial H}{\partial x}\right)^\top,\quad
\lambda(T) = \left(\frac{\partial\phi}{\partial x}\right)^\top\]
which is our Euler-Lagrange equation.

\section{Optimal Control in Continuous Time}

Through everything we have gone through already, we can set up the simplest problem to have
the state equation
\[\dot{x} = f(x,u,t),\quad x_0, T\]
for its dynamics and the performance index
\[J = \phi(x(T)) + \int_{0}^{T}L(x,u,t)dt\]
To solve the problem we look at the
\begin{enumerate}
    \item Hamiltonian: \[H(x,u,\lambda,t) = L(x,u,t) + \lambda^\top f(x,u,t)\]
    \item Optimality Condition (O.C.): \[\frac{\partial H}{\partial u} = 0\]
    \item Euler-Lagrange Equation (E.L.):
    \[\dot{\lambda} = -\left(\frac{\partial H}{\partial x}\right)^\top,
    \quad\lambda(T) = \left(\frac{\partial\phi}{\partial x}\right)^\top\]
    \item State Equation (S.E.):
    \[\dot{x} = \left(\frac{\partial H}{\partial \lambda}\right)^\top, \quad x(0) = x_0\]
\end{enumerate}
Simple Example:

\noindent
Consider the scalar linear system
\[\dot{x} = ax + bu\]
with the performance index
\[J = \frac{1}{2}\int_{0}^{T}(x^2+\rho u^2)dt\]
The Hamiltonian is
\[H = \frac{1}{2}x^2 + \frac{\rho}{2}u^2 + \lambda(ax + bu)\]
The optimality condition is
\[\frac{\partial H}{\partial u} = \rho u + b\lambda = 0\]
\[u = -\frac{b}{\rho}\lambda\]
The Euler-Lagrange equation is
\[\dot{\lambda} = -(x + a\lambda)\]
\[\lambda(T) = 0\]
Plugging everything back into our state equation:
\[\dot{x} = ax + b\left(-\frac{b}{\rho}\lambda\right)\]
We can set up an equation like so:
\[\left[\begin{matrix}
    \dot{x} \\
    \dot{\lambda}
\end{matrix}\right] = \left[\begin{matrix}
    a & -\frac{b^2}{\rho} \\[2mm]
    -1 & -a
\end{matrix}\right]\left[\begin{matrix}
    x \\ \lambda
\end{matrix}\right] = M\left[\begin{matrix}
    x \\ \lambda
\end{matrix}\right]\]
Then, we see that
\[\left[\begin{matrix}
    x(t) \\
    \lambda(t)
\end{matrix}\right] = e^{Mt}\left[\begin{matrix}
    x_0 \\ \lambda(0)
\end{matrix}\right]\]
Then,
\[\left[\begin{matrix}
    x(T) \\
    \lambda(T)
\end{matrix}\right] = \left[\begin{matrix}
    x(T) \\ 0
\end{matrix}\right] = e^{MT}\left[\begin{matrix}
    x_0 \\ \lambda(0)
\end{matrix}\right]\]
\[\left[\begin{matrix}
    x_0 \\ \lambda(0)
\end{matrix}\right] = e^{-MT}\left[\begin{matrix}
    x(T) \\ 0
\end{matrix}\right]\]
Our unknowns are $\lambda(0)$ and $x(T)$ which we will have to find. Doing this
will let us then solve for $x(t)$ and $\lambda(t)$ at any time $t$. This is the hardest part
of solving an optimal control problem.

\subsection{O.C. with (Some) Specified States at a Final Time}

Let us consider the problem where we have the dynamics
\[\dot{x} = f(x,u,t)\]
\[\dim(x) = n\]
with
\[(x_1(T),\ldots,x_p(T)),\quad p \le n\]
prescribed.
When we go through our optimal control problem and look at the perturbations to see how we
can find the optimum, we end up with
\[\delta J = \int_{0}^{T}\frac{\partial H}{\partial u}\delta u dt = 0\]
Normally, since $\delta u$ is arbitrary due to us not having specified final states,
we can iust say that $\frac{\partial H}{\partial u} = 0$. However, $\delta u$ might affect
the end states, so it can no longer be arbitrary. Therefore, we can no longer have that
optimality condition. We must compute the influence of $\delta u$ on $\delta x_i(T)$ ($i = 1,...,p$).
We specify that our overall Hamiltonian for the problem (finding extremum of performance index
$J$) is
\[H^{(J)} = L + (\lambda^{(J)})^\top f\]
Our Langrange multipliers $\lambda^{(J)}$ for this problem. We then want to
find how our perturbations affect the final states. We look at:
\[\bar{x}_i(T) = x_i(T) + \int_{0}^{T}(\lambda^{(i)})^\top(f(x,u,t)-\dot{x})dt\]
Where $\lambda^{(i)}$ is completely free. We define the Hamiltonian for this subproblem to be
\[H^{(i)} = (\lambda^{i})^\top f\]
Our choice of $\lambda^{(i)}$ satisfies its Euler-Lagrange equation:
\[\dot{\lambda}^{(i)} = -\left(\frac{\partial H^{(i)}}{\partial x}\right)^\top\]
\[\lambda^{(i)}(T) = \frac{\partial x_i(T)}{\partial x} = e_i\]
where $e_i$ is the $i$-th column of the identity matrix. We end up with the necessary condition
\[\delta x_i(T) = \int_{0}^{T}(\lambda^{(i)})^\top \frac{\partial f}{\partial u}\delta u\text{ }dt
= \int_{0}^{T}\frac{\partial H^{(i)}}{\partial u}\delta u\text{ }dt = 0\]
The trick is then to let our optimal $\delta u$ be
\[\delta u^{*}(\nu_1,\ldots,\nu_p,k) = k\left[\frac{\partial H^{(J)}}{\partial u} +
\sum_{i=1}^{p}\left(\frac{\partial H^{(i)}}{\partial u}\right)^\top \nu_i\right]\]
The new problem is now to find $\nu$ such that the final state constraints are satisfied.
\[\delta x_i(T) = k\int_{0}^{T}(\lambda^{(i)})^\top \frac{\partial f}{\partial u}
\left[\frac{\partial H^{(J)}}{\partial u} +
\sum_{i=1}^{p}\left(\frac{\partial f}{\partial u}\right)^\top\lambda^{(i)} \nu_i\right]dt\]
We set $\delta x_i(T) = 0$, which is our obiective. Therefore, $k$ is irrelevant.
We say that
\[q_i = \int_{0}^{T}(\lambda^{(i)})^\top\frac{\partial f}{\partial u}
\left(\frac{\partial H^{(J)}}{\partial u}\right)^\top dt\]
\[Q_{ii} = \int_{0}^{T}(\lambda^{(i)})^\top\left(\frac{\partial f}{\partial u}\right)
\left(\frac{\partial f}{\partial u}\right)^\top\lambda^{(i)}dt\]
So that our problem becomes
\[\delta x_i(T) = q_i + \sum_{i=1}^{p}Q_{ii}\nu_i = q_i + Q_i\nu = 0\]
Therefore, we can solve for $\nu$:
\[\nu = -Q^{-1}q\]
Note that $Q$ must be symmetric and invertible. We see that $Q$ is also the Reachability Gramian.
Basically, our problem hinges on $Q$ not having a zero eigenvalue. If $\det Q = 0$, then
there is no solution to the problem. In other words, if our system is not reachable, we are unable
to reach those final state conditions. If $Q$ satisfies these conditions, then $\delta x_i(T) = 0$
is guaranteed.

\noindent
From this, we get
\[\delta J = \int_{0}^{T}\|\delta u^*\|^2dt = 0\]
so our condition is now that $\delta u^*(\nu^*, k) = 0$

\noindent
We can produce a super Hamiltonian for the overall problem:
\begin{align*}
    H^\text{super} &= H^{(J)} + \sum_{i=1}^{p}\nu_i H^{(i)} \\
    &= L + (\lambda^{(J)})^\top f + \sum_{i=1}^{p}(\lambda^{(i)})^\top f \nu_i \\
    &= L + \left[\lambda^{(J)} + \sum_{i=1}^{p}\nu_i \lambda^{(i)}\right]^\top f
\end{align*}
So we can say we also have a super Langrange multiplier
\[\lambda^\text{super} = \lambda^{(J)} + \sum_{i=1}^{p}\nu_i \lambda^{(i)}\]
so that
\[H^\text{super} = L + \lambda^\text{super}f\]
and then our optimality condition boils down to
\[\frac{\partial H^\text{super}}{\partial u} = 0\]
Example:

\[\dot{x} = x + u, \quad x(0) = 0 \rightarrow x(1) = 1\]
Our goal is to find the minimal energy
\[J = \frac{1}{2}\int_{0}^{T}u^2 dt\]
The Hamiltonian is
\[H = \frac{1}{2}u^2 + \lambda(x + u)\]
The optimality condition is
\[\frac{\partial H}{\partial u} = 0 \rightarrow u = -\lambda\]
The Euler-Lagrange equation says
\[\dot{\lambda} = -\lambda\]
Now, how do we find the final $\lambda$? Originally, the final $\lambda$ is
\[\lambda(T) = \left(\frac{\partial\phi}{\partial x}\right)^\top\]
However, in this case, it also depends on $x$. So, we create an expression:
\[\psi = \phi + \nu x\]
and then say that
\[\lambda(T) = \left(\frac{\partial\psi}{\partial x}\right)^\top\]
In this problem, $\phi = 0$, so
\[\lambda(T) = \frac{\partial}{\partial x}(\nu x) = \nu\]
We see then that
\[\lambda(t) = e^{T-t}\nu\]
We then plug this back into our state equation:
\[\dot{x} = x + u = x - e^{T-t}\nu\]
which produces
\[x(t) = e^t x(0) - \int_{0}^{t}e^{t-\tau}e^{T-\tau}\nu\text{ }d\tau = -\nu e^T \sinh(t)\]
Then,
\[x(T) = 1 = \nu e^T \sinh(T)\]
\[\nu = -\frac{e^{-T}}{\sinh(T)}\]
So,
\[u^*(t) = -\lambda(t) = -e^{T-t}\nu = e^{T}e^{-t}\frac{e^{-T}}{\sinh(T)} = \frac{e^{-t}}{\sinh(T)}\]

\subsection{O.C. with an Unknown Final Time}

Because the time $T$ is also unknown, we need another necessary condition
to solve the problem. Consider the general Bolza problem:
\[\dot{x} = f(x,u),\quad x(0) = x_0\]
Find
\[\min_{u,T}\int_{0}^{T}L(x,u)dt + \phi(x(T),T)\]
Let $(T, x(T))$ satisfy some manifold (smooth) condition
\[\psi(T,x(T)) = 0\]
The solution to the Bolza problem then utilizes the Hamiltonian and the manifold constraint:
\[H = L + \lambda^\top f\]
\[\Psi = \phi + \nu^\top\psi\]
We introduce a new condition called the \textbf{transversality condition}
\[H + \left.\frac{\partial\Psi}{dt}\right|_{t=T} = 0\]
to help us find the final time $T$. Note that $H$ is a constant, so it does not depend on $t$.

\noindent
Example:\\[2mm]
Consider the dynamics
\[\dot{x} = u,\quad x(0) = 0\]
We want to minimize the cost
\[J = \frac{1}{2}\int_{0}^{T}(x^2+u^2)dt + \frac{\rho}{2}\left[(T-1)^2+(x-1)^2\right]\]
for some $T > 0$. Then,
\[H = \frac{1}{2}(x^2+u^2) + \lambda u\]
\[\Psi = \frac{\rho}{2}\left[(t-1)^2+(x-1)^2\right]\]
because we have a final cost $\phi$ but no manifold condition.

\noindent
Optimality Condition:
\[\frac{\partial H}{\partial u} = 0
\rightarrow u + \lambda = 0\]
\[u(t) = -\lambda(t)\]
Euler-Lagrangian:
\[\dot{\lambda} = -\left(\frac{\partial H}{\partial x}\right) = -x\]
\[\lambda(T) = \left.\frac{\partial\psi}{\partial x}\right|_{t=T} = \rho(x(T)-1)\]
State Equation:
\[\dot{x} = u = -\lambda = x(0) + \int_{0}^{t}x(\tau)d\tau\]
\[\ddot{x} = x\]
\[x(t) = A\sinh(t) + B\cosh(t)\]
Use initial condition to restrict further:
\[x(0) = A\sinh(0) + B\cosh(0) = B = 0\]
\[x(t) = A\sinh(t)\]
We see that
\[\lambda(t) = -\left[x(0) + \int_{0}^{t}x(\tau)d\tau\right]\]
\[\lambda(t) = -\int_{0}^{t}A\sinh(\tau)d\tau\]
\[\lambda(t) = -A\cosh(t)\]
At final time $T$,
\[\lambda(T) = -A\cosh(T) = \rho(A\sinh(T)-1)\]
Then,
\[A = \frac{\rho}{\rho\sinh(T)+\cosh(T)}\]
So,
\[x(t) = \frac{\rho\sinh(t)}{\rho\sinh(T)+\cosh(T)}\]
\[u(t) = -\lambda(t) = \frac{\rho\cosh(t)}{\rho\sinh(T)+\cosh(T)}\]
Transversality Condition:
\[H + \left.\frac{\partial\Psi}{\partial t}\right|_{t=T} = 0\]
\[H = \frac{1}{2}\frac{\rho^2}{(\rho\sinh(T)+\cosh(T))^2}\left(\sinh^2(t)+\cosh^2(t)\right)
+ \lambda u\]
\[= \frac{1}{2}\frac{\rho^2}{(\rho\sinh(T)+\cosh(T))^2}\left(\sinh^2(t)+\cosh^2(t)\right)
- \frac{\rho^2}{(\rho\sinh(T)+\cosh(T))^2}\cosh^2(t)\]
\[= \frac{1}{2}\frac{\rho^2}{(\rho\sinh(T)+\cosh(T))^2}\left(\sinh^2(t)-\cosh^2(t)\right)\]
\[= -\frac{\rho^2}{2(\rho\sinh(T)+\cosh(T))^2}\]
which we see is a constant.
\[\frac{\partial\Psi}{\partial t} = \rho(t-1) \rightarrow \rho(T-1)\]
Then,
\[-\frac{\rho^2}{2(\rho\sinh(T)+\cosh(T))^2} + \rho(T-1) = 0\]
\[\rho = 2(T-1)(\rho\sinh(T)+\cosh(T))^2\]
Then analyze this to see that as $\rho$ grows, our $T$ asymptotically approaches
$T = 1$ which is our minimal $T$.

\section{Optimal Control with Path Constraints}

\subsection{Integral Constraints}
Consider the Bolza problem with an additional constraint
\[I = \int_{0}^{T}N(x,u,t)dt\]
where $N$ may have many components. We then extend the state to
\[x^\text{ext} = \left[\begin{matrix}
    x \\ y
\end{matrix}\right]\]
where
\[y(t) = \int_{0}^{t}N(x,u,\tau)d\tau\]
The state equation then becomes
\[\frac{d}{dt}x^\text{ext} = \left[\begin{matrix}
    f(x,u,t) \\
    N(x,u,t)
\end{matrix}\right] = F(x^\text{ext},u,t)\]
with final state condition
\[x^\text{ext}(T) = \left[\begin{matrix}
    x(T) \\
    I
\end{matrix}\right]\]
where $x(T)$ might be unknown.
The Hamiltonian becomes
\begin{align*}
    H_\text{ext} &= L(x,u,t) + \lambda_\text{ext}^\top F(x^\text{ext},u,t) \\
    &= L(x,u,t) + \lambda^\top f(x,u,t) + \lambda_y^\top N(x,u,t) \\
    &= H_\text{old} + \lambda_y^\top N(x,u,t)
\end{align*}
So, we see we've simply just added an additional term to the original Hamiltonian and
therefore can affect the optimal control. The optimality condition becomes
\[\frac{\partial}{\partial u}H_\text{ext} + \frac{\partial}{\partial u}H_\text{old}
+ \frac{\partial}{\partial u}N(x,u,t) = 0\]
The Euler-Lagrangian produces
\[\left[\begin{matrix}
    \dot{\lambda} \\
    \dot{\lambda}_y
\end{matrix}\right] = -\frac{\partial H_\text{ext}}{\partial x^\text{ext}}
= \left[\begin{matrix}
    -\frac{\partial H_\text{ext}}{\partial x} \\[2mm]
    -\frac{\partial H_\text{ext}}{\partial y}
\end{matrix}\right] = \left[\begin{matrix}
    -\left(\frac{\partial H_\text{old}}{\partial x}+\lambda_y\frac{\partial N}{\partial x}\right) \\[2mm]
    0
\end{matrix}\right]\]
\begin{align*}
    \dot{\lambda} &= -\frac{\partial H_\text{old}}{\partial x} - \lambda_y\frac{\partial N}{\partial x} \\
    \dot{\lambda}_y &= 0
\end{align*}
with
\[\left[\begin{matrix}
    \lambda(T) \\
    \lambda_y(T)
\end{matrix}\right] = \left[\begin{matrix}
    \frac{\partial}{\partial x}\left[\phi(x)+\nu^\top(y-I)\right] \\[2mm]
    \frac{\partial}{\partial y}\left[\phi(y)+\nu^\top(y-I)\right]
\end{matrix}\right] = \left[\begin{matrix}
    \left.\frac{\partial\phi}{\partial x}\right|_{T} \\[2mm]
    \nu(T)
\end{matrix}\right]\]
We use the fact that $y(T) = I$ to solve for $\nu$.

\noindent
\underline{Example: Dido's Problem}\\[2mm]
Consider some beach where we want to claim an area of land on the beach
using just a single rope connecting two points along the shore (axis dividing ocean from land).
How do we obtain the most amount of land using a rope with a fixed length $T$?

\noindent
Let the equation modeling the position of the rope on the land $y$ with respect to the shore $x$
to be $y(x)$. No matter how the rope is laid out, we can always mark a middle point between the
two ends of the rope. Let us consider that as the origin so that the rope's ends are at points
$x = -a$ and $x = a$.

\noindent
We consider a small portion of the rope $d\ell$. $d\ell$ is made up of two component parts:
$dx$ and $dy$, which is the change in the $x$ position and the $y$ position along this portion.
This forms a right triangle, so we know that $d\ell^2 = dx^2 + dy^2$. We see then that
$dy/dx = \tan\theta$, which is the change in the rope's position depending on the point
along the shore. This is our state equation for the system.

\noindent
To set up the problem, we let
\begin{itemize}
    \item $x$: independent variable
    \item $y$: state
    \item $\theta$: control
\end{itemize}
We want to maximize the amount of area, which is modeled by
\[J = \int_{-a}^{a}y(x)dx\]
Recall, however, that the rope is a fixed length $T$. So then it must hold that
\[T = \int_{-a}^{a}d\ell(x) = \int_{-a}^{a}\sqrt{dx^2+dy^2} =
\int_{-a}^{a}\sqrt{1+\left(\frac{dy}{dx}\right)^2}dx = \int_{-a}^{a}\frac{dx}{\cos\theta}\]
which is our integral constraint. We constructed this due to the fact that $d\ell$ is
a small portion of the length of our rope, so the sum across the ends of the rope of those
small portion lengths (the integral) will produce the total length. Let
\[p(x) = \int_{-a}^{x}\frac{1}{\cos\theta}dx\]
measure the current length of the rope at point $x$ from the starting point. We produce
boundary conditions $p(-a) = 0$ and $p(a) = T$.

\noindent
The Hamiltonian is
\[H = y + \lambda\tan\theta + \lambda_p\frac{1}{\cos\theta}\]
The optimality condition is
\[\frac{\partial H}{\partial \theta} = 0 = \frac{\lambda}{\cos^2\theta}
+ \frac{\lambda_p\sin\theta}{\cos^2\theta} \rightarrow \lambda(x) = -\lambda_p\sin(\theta(x))\]
The Euler-Lagrangian yields
\begin{align*}
    \dot{\lambda} &= -\frac{\partial H}{\partial y} = -1 \\
    \dot{\lambda}_y &= 0
\end{align*}
From these differential equations, we see that
\[\lambda(x) = c - x\]
\[\lambda_p = c'\]
where $c, c'$ are constants. Plugging back into the result of our optimality condition we get
\[\lambda(x) = -\lambda_p\sin\theta\]
\[c - x = -c'\sin\theta\]
\[\sin\theta = \frac{c-x}{c'}\]
This is not enough information to plug back into our state equation just yet. Let us
make use of the fact that the Hamiltonian is constant. We see that
\[H = y - \lambda_p\sin\theta\tan\theta + \lambda_p\frac{1}{\cos\theta}
= y + \lambda_p\frac{1-\sin^2\theta}{\cos\theta} = y + \lambda_p\cos\theta\]
Then,
\[y - H  = -\lambda_p\cos\theta\]
\[x - c =  -\lambda = \lambda_p\sin\theta\]
We see that
\[(y-H)^2 + (x-c)^2 = \lambda_p^2\cos^2\theta + \lambda_p^2\sin^2\theta = \lambda_p^2\]
which eliminates $\theta$. This is a circle centered at the point $(c, H)$ with radius $\lambda_p$.
How do we solve for $H$ and $c$? Recall our boundary conditions:
\[y(-a) = 0, \quad y(a) = 0\]
\[p(-a) = 0, \quad p(a) = T\]
For the rope length, we get
\[T = \int_{-a}^{a}\frac{dx}{\cos\theta} = \int_{-a}^{a}\frac{dx}{d\theta}\frac{1}{\cos\theta}d\theta
= \int_{-a}^{a}\lambda_p d\theta = \lambda_p(\theta(a)-\theta(-a))\]
We get this because $x = c + \lambda_p\sin\theta$ so $dx/d\theta = \lambda_p\cos\theta$.
Using the boundary conditions:
\[y(-a) = 0 = H - \lambda_p\cos(\theta(-a))\]
\[y(a) = 0 = H - \lambda_p\cos(\theta(a))\]
Then,
\[H - \lambda_p\cos(\theta(-a)) = H - \lambda_p\cos(\theta(a))\]
\[\cos(\theta(-a)) = \cos(\theta(a))\]
\[-\theta(-a) = \theta(a) = \theta_0\]
since $\cos(-\theta) = -\cos\theta$. This means the angle at the start and end of the rope
are the same. Now, we solve for $c$:
\[x = -a \rightarrow c + \lambda_p\sin(\theta(-a)) = -a = c - \lambda_p\sin\theta_0\]
\[x = a \rightarrow c + \lambda_p\sin(\theta(a)) = a = c + \lambda_p\sin\theta_0\]
Then,
\[-c + \lambda_p\sin\theta_0 = c + \lambda_p\sin\theta_0\]
\[-c = c\]
\[c = 0\]
we have solved for $c$. To solve for $H$, we saw earlier that
\[y(a) = 0 = H - \lambda_p\cos\theta_0\]
So,
\[H = \lambda_p\cos\theta_0\]
We then solve for $T$:
\[T = \lambda_p(\theta_0 - (-\theta_0)) = \lambda_p 2\theta_0\]
Finally, we need to solve for $\lambda_p$. We see that
\[\lambda_p = \frac{T}{2\theta_0} \rightarrow H = \frac{T\cos\theta_0}{2\theta_0}\]
\[x = a = c + \lambda_p\sin\theta_0 = \lambda_p\sin\theta_0 \rightarrow \lambda_p
= \frac{a}{\sin\theta_0}\]
Then,
\[\frac{T}{2\theta_0} = \frac{a}{\sin\theta_0} \rightarrow
\frac{\sin\theta_0}{\theta_0} = \frac{2a}{T}\]
In order for the solution to exist, it must hold that $2a < T$.

\subsection{Equality Constraints}

\subsubsection{Equality Constraints on Control}

\noindent
Let there be some equality constraint $C(u) = 0$. We can adjoin the original problem
with a time-varying Lagrange multiplier:
\[H(x,u,\lambda,\mu) = L(x,u) + \lambda^\top f(x,u) + \mu^\top C(u)\]
Our constraint will only impact the optimality condition, forcing the control variable
to meet the equality constraint.

\noindent
\underline{Example:}\\[2mm]
Consider the problem of going from location $A$ to destination $B$ in minimum time.
We control the direction of movement (a vector) with $\hat{u}$, which is a unit vector
representing the unit direction in 3-D space.

\noindent
Our system is
\[\dot{r} = \hat{u}\]
with condition
\[\hat{u}^\top\hat{u} = 1\]
We then can say that
\[C(u) = \frac{1}{2}(\hat{u}^\top\hat{u} - 1)\]
The minimum time problem has performance index
\[J = \int_{0}^{T}dt = T\]
The Hamiltonian is
\[H = 1 + \lambda^\top u + \mu\frac{1}{2}(u^\top u - 1)\]
The optimality condition is
\[\frac{\partial H}{\partial u} = 0 = \lambda + \mu u
\rightarrow u(t) = -\frac{\lambda(t)}{\mu(t)}\]
Looking back at our condition:
\[\|u\|^2 = 1 \rightarrow \frac{\|\lambda\|^2}{\mu^2} = 1 \rightarrow \mu = \pm\|\lambda\|\]
Then,
\[u = \mp\frac{\lambda}{\|\lambda\|}\]
The Euler-Lagrangian yields
\[\dot{\lambda} = -\left(\frac{\partial H}{\partial x}\right)^\top = 0\]
which means $\lambda(t) = \text{constant}$.
The state equation is
\[\dot{r} = \hat{u} = \mp\frac{\lambda}{\|\lambda\|} = \hat{u}_0\]
\[r(t) - r(0) = (t - 0)\hat{u}_0\]
\[r(t) = r(0) + t\hat{u}_0\]
\[r(T) = r(0) + T\hat{u}_0\]
\[r_B = r_A + T\hat{u}_0\]
\[\hat{u}_0 = \frac{r_B-r_A}{T} = \frac{\lambda}{\|\lambda\|}\]
Then, $T = \|r_B-r_A\|$, so
\[u = \frac{r_B-r_A}{\|r_B-r_A\|}\]
solves our minimum time problem.

\subsubsection{Equality Constraints on State and Control}

\noindent
Let there exist some condition $C(x,u,t)$ such that $C(x,u,t) = 0, \forall t\in[0,T]$.
We see that this is a scalar value, so our Lagrange multiplier will also be a scalar.
We adjoin the constraint to the original problem as before. Our Hamiltonian becomes
\[H = L + \lambda^\top f + \mu C\]
The big difference from just a constraint on control is that the Euler-Lagrangian also changes.
We see that
\[\dot{\lambda} = -\left(\frac{\partial L}{\partial x}\right)^\top
- \left(\frac{\partial f}{\partial x}\right)^\top \lambda -
\mu\left(\frac{\partial C}{\partial x}\right)^\top\]

\noindent
\underline{Example:}\\[2mm]
Consider the dynamics
\[\dot{x} = u, \quad x(0) = x_0\]
Find for fixed $T$ the control $u$ such that
\[J = \frac{1}{2}\int_{0}^{T}\|u\|^2 dt\]
is minimal following the constraint $a^\top(x+u) = 0$.
The Hamiltonian is
\[H = \frac{1}{2}u^\top u + \lambda^\top u + \mu a^\top(x+u)\]
The optimality condition is
\[\frac{\partial H}{\partial u} = 0 = u + \lambda + \mu a \rightarrow u(t) = -\lambda(t) - \mu(t) a\]
The Euler-Lagrangian yields
\[\dot{\lambda} = -\left(\frac{\partial H}{\partial x}\right)^\top = -\mu a\]
\[\lambda(T) = 0\]
The state equation is
\[\dot{x} = u = -\lambda - \mu a\]
Plugging back into our condition:
\[a^\top x = -a^\top u = -a^\top\dot{x}\]
Then,
\[(a^\top x)(t) = e^{-t}(a^\top x)(0)\]
\[a^\top x(t) = e^{-t}a^\top x(0)\]
Looking at the Lagrangian,
\[\lambda(t) = \lambda(0) - \int_{0}^{t}\mu(\tau)d\tau\]
Then,
\[\lambda(T) = \lambda(0) - \int_{0}^{T}\mu(\tau)d\tau\]
\[\lambda(t) = \int_{t}^{T}\mu(\tau)d\tau\]
Plugging back into the control input:
\[u(t) = -a\int_{t}^{T}\mu(\tau)d\tau - a\mu(t)\]
Let
\[m(t) = \int_{0}^{t}\mu(\tau)d\tau\]
Then,
\[u(t) = am(t) - am(T) - a\mu(t)\]
\[\dot{x} = am(t) - am(T) - a\mu(t)\]
\[x(t) = x_0 - am(T)t - am(t) + a\int_{0}^{t}m(\tau)d\tau\]
\[a^\top x(t) = a^\top x_0 - a^\top am(T)t - a^\top am(t) + a^\top a\int_{0}^{t}m(\tau)d\tau\]
The left term is also
\[a^\top x(t) = e^{-t}a^\top x_0\]
Differentiate:
\[-\|a\|^2 m(T) + \|a\|^2[m(t)-\dot{\mu}(t)] = 0\]
\[m(t) - \dot{\mu}(t) = m(T)\]
\[m(t) - \ddot{m}(t) = m(T)\]
Then,
\[m(t) = A\cosh(t) + B\sinh(t) + E\]

\subsubsection{State Equality Constraints}

\noindent
Consider the dynamics
\[\dot{x} = f(x,u,t), \quad x(0) = x_0\]
Our goal is to find
\[\min_{u(\cdot)}\phi(x(T)) + \int_{0}^{T}L(x,u)dt\]
with $T$ fixed. We have a condition on the state
\[S(x(t)) = 0, \quad \forall t\]
Yet again, we adjoin it with our original problem. The Hamiltonian becomes
\[H = L + \lambda^\top f + \mu S\]
In the end, after looking at the optimality condition, Euler-Lagrangian, and state equations,
we can also plug things back into the state condition $S(x) = 0$ as a new necessary condition.
However, since $S(x)$ does not rely on $u$, it doesn't provide additional information about the
relationship between the control input and the condition. So, compute the Lie derivatives
\[L_f S, L_f^2 S, \ldots\]
until we find $u$, and then we set that to 0.

\noindent
\underline{Example:}\\[2mm]
Consider the system evolving in $\mathbb{R}^n$ with dynamics
\[\dot{x} = u, \quad \|u\| = 1\]
However, the system-state is constrained to a hyper-sphere of radius $R$
centered at the origin. Consider the states $x_0$ and $x_1$ on this sphere, as initial
and desired final states, and determine the minimum-time trajectory for this
state-constrained system.

\noindent
To set up the problem, we set up our variables:
\begin{itemize}
    \item $u$: control
    \item $x$: state
    \item $t$: independent
\end{itemize}
The equation for an n-dimensional hypersphere with radius $R$ is
\[x_1^2 + x_2^2 + \cdots + x_n^2 = R^2\]
Then a point on the hypersphere would be subject to the constraint
\[x^\top x = R^2\]
Our constraints then are
\[C(u) = u^\top u - 1 = 0\]
\[S(x) = x^\top x - R^2 = 0\]
But we see that $u$ does not appear here, so we take the Lie derivative:
\[L_f S(x) = \frac{\partial S}{\partial x}f(x,u) = 2x^\top u\]
Then, our constraint is
\[2x^\top u = 0\]
Our goal is to solve the minimum time problem,
\[J = \int_{0}^{T}dt = T\]
Our Hamiltonian is
\[H = 1 + \lambda^\top u + \mu_1(u^\top u - 1) + \mu_2(2x^\top u)\]
and terminal condition
\[\phi = \nu^\top(x-x_1)\]
The optimality condition is
\[\frac{\partial H}{\partial u} = 0 = \lambda + 2\mu_1 u + 2\mu_2 x
\rightarrow \lambda = -2\mu_1 u - 2\mu_2 x\]
The Euler-Lagrangian yields
\[\dot{\lambda} = -\left(\frac{\partial H}{\partial x}\right)^\top = -2\mu_2 u\]
\[\lambda(T) = \left(\frac{\partial\phi}{\partial x}\right)_{t=T}^\top = \nu\]
We can compare like this:
\[\dot{\lambda} = -2\dot{\mu_1}u - 2\mu_1\dot{u} - 2\dot{\mu_2}x - 2\mu_2\dot{x}
= -2\dot{\mu_1}u - 2\mu_1\dot{u} - 2\dot{\mu_2}x - 2\mu_2 u = -2\mu_2 u\]
So,
\[\dot{\mu_1}u + \mu_1\dot{u} + \dot{\mu_2}x = 0\]
\[\frac{d}{dt}(\mu_1 u) + \dot{\mu_2}x = 0\]
Multiply by $u^\top$ to get
\[u^\top\frac{d}{dt}(\mu_1 u) + \dot{\mu_2}u^\top x = 0\]
\[\dot{\mu_1} + \mu_1 u^\top\dot{u} = 0\]
We know that $u^\top\dot{u} = \frac{d}{dt}(\frac{1}{2}u^\top u) = 0$, so
\[\dot{\mu_1} = 0 \rightarrow \mu_1 = \text{constant}\]
Looking back, we see then that
\[\dot{\mu_2}x = -\dot{\mu_1}u - \mu_1\dot{u} = -\mu_1\ddot{x}\]
How do we find $\mu_2$? The optimality condition was
\[\lambda + 2\mu_1 u + 2\mu_2 x = 0\]
So then
\[\lambda^\top x + 2\mu_1 u^\top x + 2\mu_2 x^\top x = 0\]
\[\lambda^\top x + 2\mu_2 R^2 = 0\]
\[\mu_2 = -\frac{1}{2R^2}\lambda^\top x\]
Plugging back into the EL:
\[\dot{\lambda} = -2\mu_2 u = \frac{1}{R^2}\lambda x^\top u = 0\]
So $\lambda$ is a constant. Then,
\[\dot{\mu_2} = -\frac{1}{2R^2}\lambda^\top \dot{x} = -\frac{1}{2R^2}\lambda^\top u
= \frac{1}{R^2}(\mu_1 u + \mu_2 x)^\top u = \frac{1}{R^2}\mu_1 u^\top u +
\frac{1}{R^2}\mu_2 x^\top u = \frac{\mu_1}{R^2}\]
\[\mu_2(t) = \frac{\mu_1}{R^2} t + \mu_2(0)\]
So, then we get
\[\frac{\mu_1}{R^2} x = -\mu_1\ddot{x}\]
\[\ddot{x} = -\frac{1}{R^2}x\]
is our new state equation. So therefore all potential solutions are of the form
\[x(t) = \frac{A}{R^2}\cos(t) + \frac{B}{R^2}\sin(t)\]
We want to solve for $A$ and $B$ so that $x(0) = x_0$ and $x(T) = x_1$.
\[x(0) = \frac{A}{R^2}\cos(0) + \frac{B}{R^2}\sin(0) \rightarrow A = R^2 x_0\]
\[x(T) = x_0\cos(T) + \frac{B}{R^2}\sin(T) \rightarrow B = \frac{x_1-x_0\cos(T)}{\sin(T)}R^2\]
So the solution is
\[x(t) = x_0\cos(t) + \frac{x_1-x_0\cos(T)}{\sin(T)}\sin(t)
= \frac{x_1\sin(t)+x_0\sin(T)\cos(t)-x_0\cos(T)\sin(t)}{\sin(T)}\]
\[x(t) = \frac{x_1\sin(t) + x_0\sin(T-t)}{\sin(T)}\]
We want to find the minimum-time trajectory. To do this, we need to find $T$.
Let us consider the fact that our trajectory must stay on the hypersphere.
Then,
\[x^\top x = \left(\frac{x_1^\top\sin(t) + x_0^\top\sin(T-t)}{\sin(T)}\right)
\left(\frac{x_1\sin(t) + x_0\sin(T-t)}{\sin(T)}\right) = R^2\]
\[\frac{x_1^\top x_1\sin^2(t)+2x_1^\top x_0\sin(T-t)\sin(t)+x_0^\top x_0\sin^2(T-t)}
{\sin^2(T)} = R^2\]
\[R^2\sin^2(T) = x_1^\top x_1\sin^2(t)+2x_1^\top x_0\sin(T-t)\sin(t)+x_0^\top x_0\sin^2(T-t)\]
With this, we can then solve for $T$ to then get the optimal trajectory.

\subsection{Inequality Constraints}

\subsubsection{Inequality Constraints on Control}

\noindent
Consider the dynamics
\[\dot{x} = f(x,u,t), \quad x(0) = x_0\]
Our goal is to find
\[\min_{u(\cdot)}\phi(x(T)) + \int_{0}^{T}L(x,u)dt\]
for a fixed time $T$ with constraint
\[C(u,t) \le 0\]
This could be solved one of two ways. The first solution is to generalize the problem:
\[\int_{0}^{T}\delta H(u)dt = 0\]
This is a Pontryagin's maximum/minimum principle (PMP) problem. $\delta H$ is a variation of $H$
with respect to $u$ taking the constraints into account. To solve the problem we want to find
a solution
\[\min_{u\in\mathcal{U}}H(x,u,\lambda,t)\]
where $\mathcal{U}$ is the admissible set of potential solutions (i.e., ones that satisfy
the condition). This then becomes a parameter optimization problem.

\noindent
Another solution method is to adjoin the inequality constraints with Karush–Kuhn–Tucker (KKT)
multipliers.
\[H = L + \lambda^\top f + \mu^\top C\]
If $C_i < 0$, then $\mu_i = 0$ (inactive constraint).
If $C_i = 0$, then $\mu_i \ge 0$ (active constraint).
This way, in all cases, $\mu^\top C \equiv 0$.

\noindent
\underline{Example: (Scalar)}\\[2mm]
Consider
\[\dot{x} = g(t)u(t), \quad |u(t)| \le 1\]
Our goal is to find
\[\min_{u}\frac{a^2}{2}\|x(T)\|^2 + \frac{1}{2}\int_{0}^{T}u^2 dt\]
with a fixed $T$. Note that $|u(t)| \le 1$ is two constraints:
\[u-1 \le 0\]
\[-1-u \le 0\]
Then,
\[H = \frac{1}{2}u^2 + \lambda gu + \mu_1(u-1) + \mu_2(-1-u)\]
\[\phi = \frac{a^2}{2}x^2\]
The optimality condition is
\[\frac{\partial H}{\partial u} = 0 = u + g\lambda + \mu_1 - \mu_2 \rightarrow
u(t) = -g(t)\lambda - \mu_1 + \mu_2\]
The Euler-Lagrangian yields
\[\dot{\lambda} = 0\]
\[\lambda(T) = \left.\frac{\partial\phi}{\partial x}\right|_{T} = a^2 x(T)\]
So,
\[\lambda(t) = a^2 x(T)\]
\[u = -g(t)a^2 x(T) - \mu_1 - \mu_2\]
The state equation is then
\[\dot{x} = -g^2(t)a^2 x(T) - \mu_1g(t) - \mu_2g(t)\]
We then consider the inactive and active cases:
\begin{itemize}
    \item no active constraints ($\mu_1 = \mu_2 = 0$)
    \[\dot{x} = -a^2 g^2(t)x(T)\]
    \[x(t_b) - x(t_a) = -a^2x(T)\int_{t_a}^{t_b}g^2(t)dt\]
    \item $u(t) \equiv 1$ ($\mu_1 \ge 0, \mu_2 = 0$)
    \[\dot{x} = g(t)\]
    comes from $\dot{x} = g(t)u(t)$
    \item $u(t) \equiv -1$ ($\mu_1 = 0, \mu_2 \ge 0$)
    \[\dot{x} = -g(t)\]
    comes from $\dot{x} = g(t)u(t)$
\end{itemize}
We can also derive the solution using PMP.
\[H = \frac{1}{2}u^2 + \lambda gu\]
with Euler-Lagrangian
\[\dot{\lambda} = 0, \quad \lambda(T) = a^2x(T)\]
\[\lambda(t) = a^2x(T)\]
Then,
\[H = \frac{1}{2}u^2 + a^2x(T)gu\]
PMP:
\[\min_{-1 \le u \le 1}H \rightarrow u(t) = -a^2x(T)g(t)\]
Thus,
\begin{align*}
    -1 < a^2x(T)g(t) < 1 &\rightarrow u(t) = -a^2x(T)g(t) \\
    a^2x(T)g(t) > 1 &\rightarrow u(t) = -1 \\
    a^2x(T)g(t) < -1 &\rightarrow u(t) = 1
\end{align*}
With KKT conditions, we have to assume that they are differentiable. With PMP,
we do not need to make this assumption.

\noindent
\underline{Special Case:}\\[2mm]
Consider the system
\[\dot{x} = f(x) + g(x)u = Ax + bu\]
with constraint $|u| \le 1$. We want to solve the minimum time problem
\[\min_{T,u}\int_{0}^{T}dt \quad \text{such that} \quad \dot{x} = Ax + bu\]
with $x_0 \rightarrow x_f$ specified. Assume $(A,b)$ is reachable.
\[H = 1 + \lambda^\top(Ax + bu)\]
\[\psi = \nu^\top (x-x_f)\]
The Euler-Lagrangian yields
\[\dot{\lambda} = -A^\top \lambda, \quad \lambda(T) = \nu\]
Then,
\[\lambda(t) = e^{-A^\top(T-t)}\nu\]
With PMP, we look for
\[\min_{u}(\lambda^\top b)u\]
Then,
\[u = \begin{cases}
    -1, & \lambda^\top b > 0 \\
    +1, & \lambda^\top b < 0 \\
    \text{free}, & \lambda^\top b \equiv 0
\end{cases}\]
This is a bang-bang controller (input instantly flips between -1 and 1).

\noindent
Transversality Condition:
\[\left(H + \frac{\partial\Psi}{\partial t}\right)_{T} = 0\]
\[1 + \lambda^\top(T)Ax(T) + \lambda^\top(T)bu(T) = 0\]
And then the rest of the problem generally is solved through the use of geometric insight.

\noindent
\underline{Example: Double Integrator}
\[\dot{x} = v\]
\[\dot{v} = u\]
with condition $|u| \le 1$. Our goal is to find
\[\min_{|u| \le 1}T\]
with boundary conditions
\[x(0) = x_0 \rightarrow x(T) = 0\]
\[v(0) = v_0 \rightarrow v(T) = 0\]
Then,
\[H = \lambda_x v + \lambda_v u\]
\[\Psi(x,t) = t + \nu_1(x-0) + \nu_2(v-0)\]
The Euler-Lagrangian yields
\[\dot{\lambda_x} = -\left(\frac{\partial H}{\partial x}\right) = 0,
\quad \lambda_x(T) = \left.\frac{\partial\Psi}{\partial x}\right|_{T} = \nu_1\]
\[\dot{\lambda_v} = -\left(\frac{\partial H}{\partial v}\right) = -\lambda_x,
\quad \lambda_v(T) = \left.\frac{\partial\Psi}{\partial v}\right|_{T} = \nu_2\]
So,
\[\lambda_x(t) = \nu_1\]
\[\lambda_v(T)-\lambda_v(t) = -\nu_1\int_{t}^{T}dt = -\nu_1(T-t)
\rightarrow \lambda_v(t) = \nu_2 + \nu_1(T-t)\]
Transversality condition:
\[\lambda_x(T)v(T) + \lambda_v(T)u(T) + \left.\frac{\partial\Psi}{\partial t}\right|_{T} = 0\]
\[\nu_1 v(T) + \nu_2 u(T) + 1 = 0\]
Since $v(T) = 0$, we get
\[u(T)\nu_2 = -1\]
Optimality condition (PMP):
\[\min_{|u| \le 1}\lambda_v(t) u\]
We see then that the optimal solution is
\[u(t) = \begin{cases}
    -1, & \lambda_v(t) > 0 \\
    +1, & \lambda_v(t) < 0
\end{cases}\]
which is a bang-bang controller with $\lambda_v(t)$ being the \textit{switching function}.
\noindent
Consider some interval for which $u(\cdot) = -1$,
\[\dot{v} = -1 \rightarrow v(t) = T-t\]
\[\dot{x} = v \rightarrow x(t) = -\frac{1}{2}(T-t)^2\]
Then,
\[x(t) = -\frac{1}{2}v^2(t)\]
For an interval for which $u(\cdot) = 1$,
\[x(t) = \frac{1}{2}v^2(t)\]

\subsubsection{Inequalities on State and Control}

\noindent
We introduce the constraint
\[C(x,u,t) \le 0\]
First method: adjoin the constraint with a KKT multiplier:
\[H = L + \lambda^\top f + \mu^\top C\]
\[\min_{u}\int_{0}^{T}L(x,u)dt\]
\[\dot{x} = f(x,u)\]
and $\mu$ is a function of time. $C$ should be smooth (infinitely differentiable).

\noindent
The second method is to use the PMP:
\[H = L + \lambda^\top f\]
For each $x,\lambda,t$ find
\[\min_{u\in\mathcal{U}}H(x,u,\lambda,t)\]
where $\mathcal{U}$ is the admissible set of control inputs. This is valid
even if $C$ is discontinuous or not differentiable.

\subsubsection{Inequalities on State Only}

\noindent
Consider the constraint
\[S(x) \le 0, \quad \forall t\]
First, we solve the optimal control problem for $S(x) < 0$ (inactive constraint).
If the solution satisfies the condition, then we are done. Otherwise, we look at
where $S(x) \equiv 0$ in some interval. It must hold that
\[\frac{d}{dt}S(x(t)) = \frac{\partial S}{\partial x}\dot{x} \equiv 0\]
over that interval. If $u$ does not appear, we keep differentiating.

\noindent
\underline{Example:}\\[2mm]
Determine the minimal energy control for the system with dynamics
\[\dot{x} = 1 + u,\]
and initial condition $0 \le x_0 \le 1$.  The end condition at time $T = 1$ is free,
but the state is constrained by $|x| < r$ for some $1 < r$.

\noindent
The performance index for minimal energy is
\[J = \frac{1}{2}\int_{0}^{1}u^2 dt\]
The state condition is
\[S(x) = |x| - r = x^2 - r^2\]
The Hamiltonian is
\[H = \frac{1}{2}u^2 + \lambda(1 + u)\]
The Euler-Lagrangian yields
\[\dot{\lambda} = 0, \quad \lambda(T) = 0\]
Then, $\lambda = 0$. So then,
\[H = \frac{1}{2}u^2\]
Using PMP:
\[\min_{u} H \rightarrow u = 0\]
Does this work? When the constraint is inactive, the solution is
\[\frac{dx}{dt} = 1 \rightarrow x(t) = x_0 + t\]
Check $S(x(1))$:
\[S(x(1)) = (x_0 + 1)^2 - r^2 < 0\]
\[x_0 + 1 < r \rightarrow x_0 < r - 1\]
which is not always true, so this solution is not good.

\noindent
Find the control across the interval where the constraint is active:
\[L_f S = \frac{\partial S}{\partial x}\dot{x} = 2x(1+u) = 0\]
Then,
\[1 + u = 0\]
\[u = -1\]
across this interval. Analytically, we see that our control should be
to use $u = 0$ until we are about to reach $x = r$ and then switch to $u = -1$
immediately. So, we can say that the control should be
\[u(t) = \begin{cases}
    0, & 0 \le t < \tau \\
    -1, & \tau \le t \le 1
\end{cases}\]
Then, we solve for $\tau$. We should switch at the point where
\[x(\tau) = x_0 + \tau = r \rightarrow \tau = r - x_0\]
So, our optimal control is
\[u(t) = \begin{cases}
    0, & 0 \le t < r - x_0 \\
    -1, & r - x_0 \le t \le 1
\end{cases}\]

\section{Optimal Feedback Control}

\subsection{Principle of Optimality}

\noindent
A question: is it possible for optimal trajectories to intersect?

\noindent
The answer is \textbf{no}. An optimal trajectory implies uniqueness.
If two trajectories intersect, then if our starting point begins on that intersection,
we have two options to choose from for the trajectory. A well proposed
problem (that can be solved optimally) should only have one unique solution,
so this should not be able to happen.

\subsection{Hamiltonian-Jacobi-Bellman}

\noindent
Consider a problem with a terminal manifold $\psi(x,t) = 0$. Points along the terminal
manifold have ``parking fees'' (final cost) $\phi(x(T),T)$. We define the value function,
or ``cost-to-go'' to be
\[J^*(x,t) = \min_{u(\cdot)}\left\{\phi + \int Ldt\right\}\]
with boundary condition
\[J^*(x,t) = \phi(x,t) \quad \text{on} \quad \psi(x,t) = 0\]
For this, we set up a PDE known as the Hamiltonian-Jacobi-Bellman Equation:
\[-\frac{\partial J^*}{\partial t} = \min_{u}\left\{L_o + \frac{\partial J^*}{\partial x}f_o\right\}
= H^o\left(x,u,\frac{\partial J^*}{\partial x},t\right)\]
\underline{Example:}\\[2mm]
Consider the dynamics
\[\dot{x} = Fx + Gu\]
with the goal to find
\[\min \frac{1}{2}\int u^\top Pu\text{ }dt\]
with $P = P^\top > 0$. We want to go from $x_0 \rightarrow x_f$ in fixed time $T$.

\noindent
We use HJB:
\[-\frac{\partial J^*}{\partial t} = \min_{u}\left\{\frac{1}{2}u^\top Pu
+ \frac{\partial J^*}{\partial x}(Fx+Gu)\right\}\]
where
\[J^*(x,t) = \phi(x,t) + \nu^\top(x-x_f) = \nu^\top(x-x_f)\]
The optimality condition is
\[\frac{\partial H}{\partial u} = Pu + G^\top\left(\frac{\partial J^*}{\partial x}\right)^\top
= 0 \rightarrow u = -P^{-1}G\left(\frac{\partial J^*}{\partial x}\right)^\top\]
Then we can get our optimal Hamiltonian:
\[H^o = \frac{1}{2}\left(\frac{\partial J^*}{\partial x}\right)G^{-1}PP^{-1}G^\top
\left(\frac{\partial J^*}{\partial x}\right)^\top + \frac{\partial J^*}{\partial x}Fx
- \frac{\partial J^*}{\partial x}GP^{-1}G^\top\left(\frac{\partial J^*}{\partial x}\right)^\top\]
\[= -\frac{1}{2}\left(\frac{\partial J^*}{\partial x}\right)GP^{-1}
G^\top\left(\frac{\partial J^*}{\partial x}\right)^\top + \frac{\partial J^*}{\partial x}Fx\]
This is a nonlinear PDE, so we cannot find a direct solution. We want to find a solution
of the form
\[\hat{J}^*(x,t) = \frac{1}{2}x^\top S(t)x + \frac{1}{2}\nu^\top Q(t)\nu + x^\top R(t)\nu\]
where $S(t) = S^\top(t)$ and $Q(t) = Q^\top(t)$ without loss of generality.
\[-\frac{\partial\hat{J}^*}{\partial t} = -\frac{1}{2}x^\top\dot{S}x
- \frac{1}{2}\nu^\top\dot{Q}\nu - x^\top\dot{R}\nu\]
\[\left(\frac{\partial\hat{J}^*}{\partial x}\right)^\top = Sx + R\nu\]
Through some additional calculations, we derive that
\[\dot{S} = SGP^{-1}G^\top S - SF - F^\top S\]
\[\dot{Q} = R^\top GP^{-1}G^\top R\]
\[\dot{R} = (F-GP^{-1}G^\top S)^\top R\]
To solve this, we first need $S(T) = S_f$ to be known. Then, we can solve for $S$,
plug that result in to find $R$, and then solve for $Q$ using that result.

\section{Linear Systems with Quadratic Performance Index}

\subsection{Terminal Controllers with Soft Constraints}

\noindent
Consider the system
\[\dot{x}(t) = F(t)x(t) + G(t)u(t), \quad x(t_0) = x_0\]
Our goal is to drive the state to the equilibrium $x = 0$ in finite time,
but without too much cost. This problem has the performance index
\[J = \frac{1}{2}x^\top(T)S_f x(T) + \frac{1}{2}\int_{0}^{T}x^\top Ax + u^\top Bu\text{ d}t\]
with $A \ge 0, B > 0, S_f \ge 0$.
Even if the system $(F,G)$ were reachable, it may not be best to reach the equilibrium
exactly due to high costs, so we tolerate a neighborhood of $x = 0$ through the parking fee
and rely on a soft constraint of $x(T) \approx 0$.

\noindent
The Hamiltonian for the optimal control problem is
\[H = \frac{1}{2}x^\top Ax + \frac{1}{2}u^\top Bu + \lambda^\top(Fx+Gu)\]
\underline{O.C.}
\[\frac{\partial H}{\partial u} = Bu + G\top\lambda = 0 \rightarrow u = -B^{-1}G^\top\lambda\]
\underline{E.L.}
\[\dot{\lambda} = -\left(\frac{\partial H}{\partial x}\right)^\top = -Ax-F^\top\lambda\]
\underline{S.E.}
\[\dot{x} = Fx - GB^{-1}G^\top\lambda\]
\underline{B.C.}
\[x(0) = x_0\]
\[\lambda(T) = \left.\left(\frac{\partial\phi}{\partial x}\right)^\top\right|_{T} = S_f x(T)\]
As you can see, this leaves us with two coupled PDEs for $x$ and $\lambda$. We have solved coupled
PDEs like this before, by putting it into a matrix form:
\[\left[\begin{matrix}
    \dot{x} \\ \dot{\lambda}
\end{matrix}\right] = \left[\begin{matrix}
    F & -GB^{-1}G^\top \\
    -A & -F^\top
\end{matrix}\right]\left[\begin{matrix}
    x \\ \lambda
\end{matrix}\right]\]
We call this the Hamiltonian system. If we assume that $\lambda(0)$ is known, then
we can write the solution as a transition function matrix:
\[\left[\begin{matrix}
    x(t) \\ \lambda(t)
\end{matrix}\right] = \Psi(t, 0)\left[\begin{matrix}
    x(0) \\ \lambda(0)
\end{matrix}\right]\]
with
\[\Psi(t, 0) = \left[\begin{matrix}
    \psi_{11} & \psi_{12} \\
    \psi_{21} & \psi_{22}
\end{matrix}\right]\]
Then,
\[x(T) = \psi_{11}(T)x(0) + \psi_{12}(T)\lambda(0)\]
\[\lambda(T) = \psi_{21}(T)x(0) + \psi_{22}(T)\lambda(0) \implies
\lambda(0) = \psi_{22}^{-1}(T)\left[S_f x(T) - \psi_{21}(T)x_0\right]\]
\[x(T) = \psi_{11}(T)x(0) + \psi_{12}(T)\psi_{22}^{-1}(T)\left[S_f x(T) - \psi_{21}(T)x_0\right]\]
\[\left[I - \psi_{12}(T)\psi_{22}^{-1}(T)S_f\right]x(T) = \left[\psi_{11}(T)
- \psi_{12}(T)\psi_{22}^{-1}(T)\psi_{21}(T)\right]x_0\]
is our solution for the optimal trajectory under these conditions. But what if
we do not know $\lambda(0)$, such is the case for almost all optimal control problems?

\subsubsection{Riccati Equation}

\noindent
Since $\lambda(T) = S_f x(T)$, let $\lambda(t) = S(t)x(t)$. Then,
\[\dot{\lambda}(t) = \dot{S}(t)x(t) + S(t)\dot{x}(t)\]
Plugging in our Hamiltonian system, we get:
\[-Ax - F^\top\lambda = \dot{S}x + S(Fx-GB^{-1}G^\top\lambda)\]
Since $\lambda = Sx$:
\[-Ax - F^\top Sx = \dot{S}x + S(Fx-GB^{-1}G^\top Sx)\]
\[\left[-A-F^\top S-\dot{S}-SF+SGB^{-1}G^\top S\right]x = 0\]
Since $x$ is arbitrary, we cannot rely on $x = 0$ as the enforcement to this rule.
So,
\[-A-F^\top S-\dot{S}-SF+SGB^{-1}G^\top S = 0\]
\[\dot{S} = -A-F^\top S+SGB^{-1}G^\top S\]
is a PDE for finding a valid $S(t)$ such that $S(T) = S_f$.
This is known as the Riccati Equation. Under this, the closed-loop state equation is
\[\dot{x} = Fx - GB^{-1}G^\top Sx = \left[F-GB^{-1}G^\top S\right]x\]
Consider
\[\int_{0}^{T}\frac{d}{dt}\left[x^\top(t)S(t)x(t)\right]dt = x^\top(T)S(T)x(T)-x^\top(0)S(0)x(0)\]
\[= \int_{0}^{T}\left[\dot{x}^\top Sx + x^\top\dot{S}x + x^\top S\dot{x}\right]dt\]
The first term yields
\[\int_{0}^{T}\dot{x}^\top Sx\text{ d}t = \int_{0}^{T}x^\top(F^\top SGB^{-1}G^\top)Sx\text{ d}t\]
The second term yields
\[\int_{0}^{T}\dot{x}^\top \dot{S}x\text{ d}t =
\int_{0}^{T}x^\top[-A-F^\top S-SF+SGB^{-1}G^\top S]x\text{ d}t\]
The third term yields
\[\int_{0}^{T}x^\top S\dot{x}\text{ d}t = \int_{0}^{T}x^\top S(F-GB^{-1}G^\top S)x\text{ d}t\]
The sum of these terms results in
\[\int_{0}^{T}x^\top(-SGB^{-1}G^\top S-A)x\text{ d}t = x^\top(T)S_fx(T)-x_0^\top S(0)x_0\]
\[= -\int_{0}^{T}x^\top Ax\text{ d}t - \int_{0}^{T}x^\top SGB^{-1}BB^{-1}G^\top Sx\text{ d}t\]
\[= -\int_{0}^{T}x^\top Ax + u^\top Bu\text{ d}t\]
So,
\[J^* = x^\top(T)S_fx(T) + \int_{0}^{T}(x^\top Ax + u^\top Bu)\text{ d}t
= x_0^\top S(0)x_0\]
This method is also known as the \textbf{sweep method}.

\noindent
\underline{Example:}

\noindent
Given the dynamics
\[\dot{x} = u, \quad x(0) = x_0\]
compute the optimal trajectory $x(t)$ for the optimal control $u(t)$ that minimizes
\[\frac{1}{2}cx^2(T) + \frac{1}{2}\int_{0}^{T}u^2(t)\text{ d}t\]
Looking at this, we find that
\[F = 0, A = 0, S_f = c, G = 1, B = 1\]
The Riccati equation therefore yields
\[\dot{S} = -A-F^\top S+SGB^{-1}G^\top S = S^2\]
Then,
\[\frac{dS}{S^2} = dt \rightarrow -\frac{1}{S(t)} + \frac{1}{S(T)} = t-T\]
\[S(t) = \frac{1}{\frac{1}{c}+T-t}\]
The feedback law is
\[u(t) = -B^{-1}G^\top S(t)x(t) = -\frac{1}{\frac{1}{c}+T-t}x(t)\]
Then, the closed-loop state equation becomes
\[\dot{x}(t) = \frac{1}{t-T-\frac{1}{c}}x(t)\]
And we get the optimal trajectory:
\[\frac{dx}{x} = \frac{d(t-T-\frac{1}{c})}{t-T-\frac{1}{c}}\]
\[\ln x(t) - \ln x_0 = \ln\left(t-T-\frac{1}{c}\right)-\ln\left(-T-\frac{1}{c}\right)\]
But, this does not work because of negatives inside the logarithm. So, we multiply
the numerator and denominator by $-1$:
\[\frac{dx}{x} = \frac{d(T-t+\frac{1}{c})}{T-t+\frac{1}{c}}\]
\[\ln x(t) - \ln x_0 = \ln\left(T-t+\frac{1}{c}\right)-\ln\left(T+\frac{1}{c}\right)\]
\[x(t) = \frac{T-t+\frac{1}{c}}{T+\frac{1}{c}}x_0\]
Our goal, remember, is to drive $x(T)$ to a neighborhood around $x = 0$ in finite time $T$.
We see that as $c \rightarrow \infty$, $x(T) \rightarrow 0$, which is the expected behavior.

\subsection{Terminal Controllers with Hard Constraints}

Consider again the same system but this time it is mandatory to reach the (linear) terminal manifold
$\psi = Mx(T)$ at a fixed final time $T$. Assume that $\psi$ has $p$ components (i.e., $M$ is
a $p \times n$ matrix).
\[\phi = \nu^\top\psi = \nu^\top Mx\]
\[H = \frac{1}{2}(x^\top Ax + u^\top Bu) + \lambda^\top(Fx+Gu)\]
\underline{O.C.}
\[\frac{\partial H}{\partial u} = 0 = Bu + G^\top\lambda \rightarrow u = -B^{-1}G^\top\lambda\]
\underline{E.L.}
\[\dot{\lambda} = -Ax - F^\top\lambda\]
\underline{S.E.}
\[\dot{x} = Fx-GB^{-1}G^\top\lambda\]
And then we end up with the Hamiltonian system:
\[\left[\begin{matrix}
    \dot{x} \\ \dot{\lambda}
\end{matrix}\right] = \left[\begin{matrix}
    F & -GB^{-1}G^\top \\
    -A & -F^\top
\end{matrix}\right]\left[\begin{matrix}
    x \\ \lambda
\end{matrix}\right]\]
There are $2n$ ODEs and $2n$ boundary conditions (initial state and terminal costate).
However, there are at least $p$ unknowns from the costates due to the terminal manifold condition
$(\nu_1,\ldots,\nu_p)$. How do we solve this problem? There are two methods:
the transition matrix method if $n-p$ terimanl variables are known or the sweep method.
As with the soft constraint version, the sweep method leads to a cleaner analytical approach,
so we will postulate a solution of the form:
\begin{align*}
    \lambda(t) &= S(t)x(t) + R(t)\nu \\
    \psi &= U(t)x(t) + Q(t)\nu
\end{align*}
with $U(T) = M, Q(T) = 0, R(T) = M^\top, S(T) = 0$ for terminal conditions.
Then,
\[\lambda = Sx + R\nu \rightarrow \dot{\lambda} = \dot{S}x + S\dot{x} + \dot{R}\nu\]
Since $\dot{x} = Fx-GB^{-1}G^\top\lambda$, we subtitute:
\[-Ax-F^\top(Sx+R\nu) = \dot{S}x + S(Fx-GB^{-1}G^\top(Sx+R\nu)) + \dot{R}\nu\]
Decouple $x$ and $\nu$:
\[(-A-F^\top S-\dot{S}-SF+SGB^{-1}G^\top S)x = (F^\top R-SGB^{-1}G^\top R+\dot{R})\nu\]
We find $\dot{S}$ and $\dot{R}$ such that both sides are rendered to 0 regardless
of any arbitrary $x$ or $\nu$. This yields:
\begin{align*}
    \dot{S} &= -A-F^\top S - SF + SGB^{-1}G^\top S \quad\leftarrow\text{Riccati Equation} \\
    \dot{R} &= -(F^\top - SGB^{-1}G^\top)R
\end{align*}
Differentiation of the $\psi$-equation yields
\[0 = \dot{U}x + U\dot{x} + \dot{Q}\nu\]
\[0 = \dot{U}x + U(F-GB^{-1}G^\top S)x - UGB^{-1}G^\top R\nu + \dot{Q}\nu\]
\[[\dot{U}+U(F-GB^{-1}G^\top S)]x + [\dot{Q}-UGB^{-1}G^\top R]\nu = 0\]
Then, to render this to zero regardless of $x$ and $\nu$,
we choose $\dot{Q}$ and $\dot{U}$ such that
\begin{align*}
    \dot{U} &= -U(F-GB^{-1}G^\top S) \\
    \dot{Q} &= UGB^{-1}G^\top R
\end{align*}
We see from earlier, then, that $U(t) = R^\top(t)$. So,
\[\dot{Q} = R^\top GB^{-1}G^\top R\]
with $Q(T) = 0$ is a reachability Gramian. Then, through the sweep method, we solve
for $S$ which yields then the solution to $R$ which yields then the solution to $Q$.
Then, we find $\nu$:
\[\nu = Q^{-1}(t)[\psi-R^\top(t)x(t)]\]
Plug this into the state equation:
\[\dot{x} = (F-GB^{-1}G^\top S)x - GB^{-1}G^\top RQ^{-1}[\psi-R^\top x]\]
This is a sort of correction term on the manifold, that considers when
$\nu = 0$. We can say that this is $\hat{\psi}$. So,
\[\dot{x} = (F-GB^{-1}G^\top S)x - GB^{-1}G^\top RQ^{-1}(\psi-\hat{\psi})\]

\subsection{Regulators}

\noindent
If $F,G,A,B$ are constant matrices (time-invariant system), then $S(t)$ may approach
a constant matrix $S_\infty$. The feedback control would be
\[u(t) = -C(t)x(t)\]
where $C \rightarrow C_\infty = B^{-1}G^\top S_\infty$.

\noindent
\underline{Example:}

\noindent
Consider the spring oscillator system:
\begin{align*}
    \dot{x}_1 &= x_2 \\
    \dot{x}_2 &= -\omega_0^2 x_1 + u
\end{align*}
The performance index is
\[J = \frac{1}{2}\int_{0}^{\infty}(ax_1^2+u^2)\text{ d}t\]
Then,
\[A = \left[\begin{matrix}
    a & 0 \\
    0 & 0
\end{matrix}\right], \quad B = 1\]
\[F = \left[\begin{matrix}
    0 & 1 \\
    -\omega_0^2 & 0
\end{matrix}\right], \quad G = \left[\begin{matrix}
    0 \\ 1
\end{matrix}\right]\]
We consider the solution to the Riccati equation $S(t)$ to be a symmetric matrix
\[S(t) = \left[\begin{matrix}
    S_1 & S_x \\
    S_x & S_2
\end{matrix}\right]\]
which satisfies
\[\dot{S} = -A - F^\top S - SF + SGB^{-1}G^\top S\]
\[\left[\begin{matrix}
    \dot{S}_1 & \dot{S}_x \\
    \dot{S}_x & \dot{S}_2
\end{matrix}\right] = -\left[\begin{matrix}
    a & 0 \\
    0 & 0
\end{matrix}\right] - \left[\begin{matrix}
    0 & -\omega_0^2 \\
    1 & 0
\end{matrix}\right]\left[\begin{matrix}
    S_1 & S_x \\
    S_x & S_2
\end{matrix}\right] - \left[\begin{matrix}
    S_1 & S_x \\
    S_x & S_2
\end{matrix}\right]\left[\begin{matrix}
    0 & 1 \\
    -\omega_0^2 & 0
\end{matrix}\right] + \left[\begin{matrix}
    0 & S_x \\
    0 & S_2
\end{matrix}\right]\left[\begin{matrix}
    S_1 & S_x \\
    S_x & S_2
\end{matrix}\right]\]
\[= -\left[\begin{matrix}
    a & 0 \\
    0 & 0
\end{matrix}\right] - \left[\begin{matrix}
    -S_x\omega_0^2 & -S_2\omega_0^2 \\
    S_1 & S_x
\end{matrix}\right] - \left[\begin{matrix}
    -S_x\omega_0^2 & S_1 \\
    -S_2\omega_0^2 & S_x
\end{matrix}\right] + \left[\begin{matrix}
    S_x^2 & S_xS_2 \\
    S_xS_2 & S_2^2
\end{matrix}\right]\]
\[= \left[\begin{matrix}
    -a+2S_x\omega_0^2+S_x^2 & -S_1+S_2\omega_0^2+S_xS_2 \\
    -S_1+S_2\omega_0^2+S_xS_2 & -2S_x+S_2^2
\end{matrix}\right]\]
\begin{align*}
    \dot{S}_1 &= -a+2S_x\omega_0^2+S_x^2 \\
    \dot{S}_x &= -S_1+S_2\omega_0^2+S_xS_2 \\
    \dot{S}_2 &= -2S_x+S_2^2
\end{align*}
The solution to the algebraic Reiccati equation (ARE) assumes that $S$ is a constant
(as posed by our problem), and therefore $\dot{S} = 0$ at the steady state (when $S = S_\infty$).
Using this, we can solve for $S_\infty$ ($S$ at steady state):
\begin{align*}
    S_1 &= \sqrt{\omega_0^4+a} \cdot \sqrt{2\sqrt{\omega_0^4+a}-2\omega_0^2} \\
    S_x &= \sqrt{\omega_0^4+a}-\omega_0^2 \\
    S_2 &= \sqrt{2\sqrt{\omega_0^4+a}-2\omega_0^2}
\end{align*}
Then, the optimal control is
\[u = -Cx\]
with
\[C = \left[\begin{matrix}
    0 & 1
\end{matrix}\right]\left[\begin{matrix}
    \sqrt{\omega_0^4+a} \cdot \sqrt{2\sqrt{\omega_0^4+a}-2\omega_0^2} & \sqrt{\omega_0^4+a}-\omega_0^2 \\
    \sqrt{\omega_0^4+a}-\omega_0^2 & \sqrt{2\sqrt{\omega_0^4+a}-2\omega_0^2}
\end{matrix}\right]\]
\[= \left[\begin{matrix}
    \sqrt{\omega_0^4+a}-\omega_0^2 & \sqrt{2\sqrt{\omega_0^4+a}-2\omega_0^2}
\end{matrix}\right] = \left[\begin{matrix}
    C_1 & C_2
\end{matrix}\right]\]
so then
\[u = -C_1x_1 - C_2x_2\]
where $C_1$ and $C_2$ are functions of $a$. The closed-loop system is
\[\ddot{x} = -\omega_0^2x_1-C_1x_1-C_2\dot{x}_1\]
Let's find the root locus of the system as a function of $a$:
\[s^2X(s) = -\omega_0^2X(s)-\left(\sqrt{\omega_0^4+a}-\omega_0^2\right)X(s)
- sX(s)\sqrt{2\sqrt{\omega_0^4+a}-2\omega_0^2}\]
\[s^2 = -\sqrt{\omega_0^4+a} - s\sqrt{2\sqrt{\omega_0^4+a}-2\omega_0^2}\]
\[\left(s^2+\sqrt{\omega_0^4+a}\right)^2=2s^2\left[\sqrt{\omega_0^4+a}-\omega_0^2\right]\]
\[s^4+2s^2\sqrt{\omega_0^4+a}+\omega_0^4+a = 2s^2\sqrt{\omega_0^4+a}-2s^2\omega_0^2\]
\[s^4+2s^2\omega_0^2+\omega_0^4+a = 0\]
\[-\frac{1}{a} = \frac{1}{(s^2+\omega_0^2)^2}\]
So, for the open loop, we have double poles at $\pm j\omega_0$ with zeros going off to infinity.
\begin{theorem}
    If the ARE has a positive semi-definite solution, then there exists a unique maximal
    solution $\tilde{S}$ of the ARE and the control
    \[u = -B^{-1}G^\top\tilde{S}x\]
    minimizes the performance index $J$ to the value function $\frac{1}{2}x_0^\top\tilde{S}x_0$.
\end{theorem}
Look back at the Hamiltonian system:
\[\left[\begin{matrix}
    \dot{x} \\ \dot{\lambda}
\end{matrix}\right] = \left[\begin{matrix}
    F & -GB^{-1}G^\top \\
    -A & -F^\top
\end{matrix}\right]\left[\begin{matrix}
    x \\ \lambda
\end{matrix}\right]\]
Let
\[\mathcal{H} = \left[\begin{matrix}
    F & -GB^{-1}G^\top \\
    -A & -F^\top
\end{matrix}\right], \quad J = \left[\begin{matrix}
    0 & I \\
    -I & 0
\end{matrix}\right]\]
We see that
\[J\mathcal{H}^\top J^\top = -\mathcal{H}\]
which makes $\mathcal{H}$ a symplectic matrix. What is special about this is that
if $s$ is an eigenvalue of $\mathcal{H}$, then so is $-s, \bar{s}, -\bar{s}$.
\begin{theorem}
    Stabilizability and Detectability
    \begin{enumerate}
        \item If $(F,G)$ is stabilzable, then the ARE has at least one solution $S = S^\top \ge 0$.
        A system is stable if there exists a feedback control law $u = -Cx$ such that the closed-loop
        system is stable, including any uncontrollable modes. If a mode is unstable, it must
        also be controllable.
        \item If $A = H^\top H$ and $(F,G)$ is detectable, then the ARE has at most one solution
        $S = S^\top \ge 0$. If $S$ is the solution, then the closed-loop dynamic matrix
        $F-GB^{-1}G^\top S = F-GC$ is stable.
    \end{enumerate}
\end{theorem}

\section{Minimum Integral Square Error}

\subsection{Minimum Integral-Square Error for SISO Systems}

A useful performance specification for the response to impulse disturbances
is to minimize the integral-square error for a specified amount of integral-square control, i.e.
\[\min_{u}\int_{0}^{\infty}y(t)^2\text{ d}t\]
with
\[\int_{0}^{\infty}u(t)^2\text{ d}t\]
specified. The open-loop transfer function is then
\[H(s) = \frac{\hat{y}(s)}{\hat{u}(s)}\]
This is equivalent to minimizing a weighted sum of the two integrals:
\[\min_{u}\int_{0}^{\infty}[Ay(t)^2+u(t)^2]\text{ d}t\]
Using the calculus of variations, we use the control
\[u(t) = -kx(t)\]
where $k$ is some constant. The feedback system should have its eigenvalues (roots)
in the left-half plane. This can be enforced by finding the root-square characteristic
equation (RSCE):
\[H(-s)AH(s) + 1 = 0.\]
The gain vector, $k$, can be obtained by equating like powers of $s^2$ in the identity
\[H(-s)AH(s) + 1 = [kX(-s)+1][kX(s)+1],\]
where
\[\frac{\hat{x}(s)}{\hat{u}(s)} = X(s), \quad H(s) = cX(s),\]
and $kX(s)+1=0$ is the closed-loop characteristic equation of the system for the given control.
We define
\[H(s) = \frac{N(s)}{D(s)}\]
Then, the root-square locus is
\[\frac{N(s)}{D(s)}\cdot\frac{N(-s)}{D(-s)} = -\frac{1}{A}\]
\underline{Example:}

\noindent
Consider the scalar system
\[\dot{y} = -by + u \implies \frac{\hat{y}(s)}{\hat{u}(s)} = \frac{1}{s+b}\]
Then the root-square locus is
\[\frac{1}{s+b}\cdot\frac{1}{-s+b} = -\frac{1}{A}\]
Then,
\[A - (s^2-b^2) = 0 \implies s = \pm\sqrt{b^2+A}\]
We see that as $A$ increases, the poles move further away from the imaginary axis.

\noindent
The root-square locus reflects poles and zeros across the imaginary axis, leading
to double the number of poles and zeros. We select the $0^\circ$ or $180^\circ$
root locus depending on whichever one has no segments crossing over the imaginary axis.
In this case, we select the $0^\circ$ root locus since it has pole starting at $-b$
and going to $-\infty$ as $A$ increases, thus making the system stable.

\noindent
We note that $A = s^2 - b^2$. The RSCE is factored to give
\[H(-s)AH(s) + 1 = \frac{A}{-s^2+b^2} + 1 = 0\]
\[\frac{A}{-s^2+b^2} + 1 \equiv \left(kX(-s) + 1\right)\left(kX(s) + 1\right)\]
We will simply let $X(s) = H(s) = \frac{1}{s+b}$:
\[\left(\frac{k}{-s+b}+1\right)\left(\frac{k}{s+b}+1\right) = \frac{A}{-s^2+b^2} + 1\]
Solving through this (via quadratic formula), we end up with
\[k = -b \pm \sqrt{b^2+A}\]
so we can simply select $k = -b + \sqrt{b^2+A}$ as our optimal feedback gain.

\noindent
\underline{Example 2:}

\noindent
Consider the undamped harmonic oscillator:
\begin{align*}
    \dot{x}_1 &= x_2 \\
    \dot{x}_2 &= -\omega^2 x_1 + \omega^2 u \\
    y &= x_1
\end{align*}
Then, the open-loop transfer function is
\[\hat{y}(s) = \hat{x}_1(s)\]
\[s\hat{x}_1(s) = \hat{x}_2(s)\]
\[s\hat{x}_2(s) = -\omega^2\hat{x}_1(s)+\omega^2\hat{u}(s)
\implies \hat{x}_2(s) = \frac{\omega^2}{s}(-\hat{x}_1(s)+\hat{u}(s))\]
\[s^2\hat{x}_1(s) = -\omega^2\hat{x}_1(s)+\omega^2\hat{u}(s)
\implies \hat{x}_1(s) = \frac{\omega^2}{s^2+\omega^2}\hat{u}(s) = \hat{y}(s)\]
\[\frac{\hat{y}(s)}{\hat{u}(s)} = \frac{\omega^2}{s^2+\omega^2}\]
The RSCE is:
\[\frac{\omega^2}{(-s)^2+\omega^2}\cdot\frac{\omega^2}{s^2+\omega^2} = -\frac{1}{A}\]
\[-\frac{1}{A} = \frac{\omega^4}{(s^2+\omega^2)^2}\]
In Evans's form, the RSCE is
\[-A\omega^4 = (s^2+\omega^2)^2\]
The roots are
\[s^2+\omega^2 = j\omega^2\sqrt{A}\]
\[s = \pm jw\sqrt{1-j\sqrt{A}}\]
The RSCE could be factored, but we can also form it through state feedback:
\[u = -k_1x_1 - k_2x_2\]
The characteristic equation of the closed-loop system with the above feedback is:
\[\left|\begin{matrix}
    s & -1 \\
    \omega^2(1+k_1) & s+\omega^2k_2
\end{matrix}\right| = 0\]
\[s^2 + k_2\omega^2s + \omega^2(1+k_1) = 0\]
Thus, the RSCE is
\[[(-s)^2+k_2\omega^2(-s)+\omega^2(1+k_1)][s^2+k_2\omega^2s+\omega^2(1+k_1)] = 0\]
Through distributing and finding the polynomial for the RSCE and then equating
like coefficients to the result of the Evans's form, we obtain the optimal feedback gains
\begin{align*}
    k_1 &= \sqrt{1+A}-1 \\
    k_2 &= \frac{\sqrt{2k_1}}{\omega}
\end{align*}

\end{document}
